\documentclass[11pt,letterpaper]{article}
\usepackage[utf8]{inputenc}
\usepackage[T1]{fontenc}
\usepackage{lmodern}
\usepackage[spanish,es-nodecimaldot]{babel}
\usepackage[margin=2.2cm]{geometry}
\usepackage{graphicx,float,booktabs,longtable,array,enumitem,amssymb}
\usepackage[hidelinks]{hyperref}
\usepackage{xcolor,fancyhdr}
% Permite compilar desde la carpeta del documento, entregables o el proyecto.
\graphicspath{{../img/}{img/}{entregables/img/}{../../img/}{../../entregables/img/}{../entregables/img/}}
\setlength{\parskip}{4pt}\setlength{\parindent}{0pt}
\pagestyle{fancy}\fancyhf{}
\fancyhead[L]{MINE-4214 · Proyecto Entrega 1}\fancyhead[R]{CupiTaller · Universidad de los Andes}
\fancyfoot[C]{\thepage}
\setcounter{secnumdepth}{2}
\begin{document}
\begin{center}
{\Large\bfseries Anexos del informe\\[2pt] Proyecto · Entrega 1 · CupiTaller}\\[8pt]
MINE-4214 Modelado y Diseño de Datos\\
Universidad de los Andes · Octubre de 2026\\[6pt]
\textbf{Integrantes:} \emph{María Alejandra Pérez Petro}, \emph{Santiago Martinez Novoa}
\end{center}
\tableofcontents
\vspace{4pt}
Anexos del informe principal (\texttt{informe\_latex.tex}). Las referencias a secciones numeradas (por ejemplo, 3.6) remiten a ese informe, y las reglas de negocio R01 a R04 se definen en su sección 1.3.
\bigskip
\section*{Anexo A · Glosario}

\begingroup\small\setlength{\tabcolsep}{4pt}\renewcommand{\arraystretch}{1.15}
\begin{longtable}{>{\raggedright\arraybackslash}p{\dimexpr 0.408\textwidth-2\tabcolsep\relax}>{\raggedright\arraybackslash}p{\dimexpr 0.522\textwidth-2\tabcolsep\relax}}
\toprule
\textbf{Término} & \textbf{Significado} \\
\midrule
\endhead
\textbf{Bookeau} & Plataforma web donde los estudiantes reservan las tutorías. \\
\textbf{Monitor} & Estudiante que da tutorías a sus compañeros. En los datos aparece como «prestador del servicio». \\
\textbf{Modalidad} & Tipo de tutoría: Normal, Express o Normal pico (ver 1.1). \\
\textbf{Servicio} & Curso para el cual se pide la tutoría (IP, EDA, APO 1, etc.). \\
\textbf{Período} & Semestre académico, con código de seis cifras (por ejemplo, 202610 es el primer semestre de 2026; terminación 19 es el intersemestral). \\
\textbf{En cola} & Solicitud que no consiguió cupo y quedó en lista de espera. \\
\textbf{Sanción} & Penalización por incumplir las reglas de reserva. \\
\textbf{Variable} & Una columna de una tabla de datos. \\
\bottomrule
\end{longtable}
\endgroup


\section*{Anexo B · Reproducibilidad}

Todo el proceso se puede volver a ejecutar desde los archivos originales con programas en Python:
\begin{enumerate}[leftmargin=1.5em,itemsep=2pt]
\item \textbf{Perfilamiento de Bronze:} lee los 114 archivos de Excel sin modificarlos y calcula estadísticas, vínculos y reglas de validez.
\item \textbf{Limpieza (Bronze a Silver):} genera los cinco archivos CSV con las reglas de la sección 3.6 y las banderas de calidad.
\item \textbf{Carga (Silver a Gold):} reconstruye la base de datos SQLite con el modelo dimensional y ejecuta los controles de la sección 5.3.
\item \textbf{Tablero:} genera el archivo HTML a partir de la base Gold.
\end{enumerate}

Cada paso reemplaza sus resultados solo cuando termina sin errores. Los resultados describen esta extracción completa de datos, no la calidad garantizada de futuras cargas ni un diccionario confirmado por el responsable de los datos. Si cambian las fuentes o las reglas de negocio, hay que volver a ejecutar los pasos y revisar las definiciones de los análisis.

\section*{Anexo C · Datos de apoyo detallados}

Este anexo reúne el detalle completo de la caracterización, la calidad, la limpieza y la carga. El informe usa un resumen de estas tablas. Las cifras describen todos los datos recibidos, antes de aplicar la regla R03 (solo reservas de tutorías), y deben recalcularse con esa población.

\subsection*{C.1 Resumen de las cinco fuentes}

\begingroup\footnotesize\setlength{\tabcolsep}{4pt}\renewcommand{\arraystretch}{1.15}
\begin{longtable}{>{\raggedright\arraybackslash}p{\dimexpr 0.194\textwidth-2\tabcolsep\relax}>{\raggedleft\arraybackslash}p{\dimexpr 0.076\textwidth-2\tabcolsep\relax}>{\raggedleft\arraybackslash}p{\dimexpr 0.057\textwidth-2\tabcolsep\relax}>{\raggedleft\arraybackslash}p{\dimexpr 0.082\textwidth-2\tabcolsep\relax}>{\raggedleft\arraybackslash}p{\dimexpr 0.088\textwidth-2\tabcolsep\relax}>{\raggedleft\arraybackslash}p{\dimexpr 0.082\textwidth-2\tabcolsep\relax}>{\raggedright\arraybackslash}p{\dimexpr 0.076\textwidth-2\tabcolsep\relax}>{\raggedright\arraybackslash}p{\dimexpr 0.076\textwidth-2\tabcolsep\relax}>{\raggedleft\arraybackslash}p{\dimexpr 0.076\textwidth-2\tabcolsep\relax}>{\raggedleft\arraybackslash}p{\dimexpr 0.070\textwidth-2\tabcolsep\relax}>{\raggedleft\arraybackslash}p{\dimexpr 0.051\textwidth-2\tabcolsep\relax}}
\toprule
\textbf{Fuente} & \textbf{Archivos} & \textbf{Filas} & \textbf{Variables} & \textbf{Variantes de encabezado} & \textbf{ID distintos} & \textbf{Primer inicio} & \textbf{Último inicio} & \textbf{Duración mín. (min)} & \textbf{Mediana} & \textbf{Máx.} \\
\midrule
\endhead
Encuesta de satisfacción Express & 29 & 6.759 & 33 & 1 & 6.759 & 2017-02-02 & 2026-05-22 & 20 & 30 & 50 \\
Encuesta de satisfacción Normal & 29 & 16.379 & 35 & 1 & 16.379 & 2016-08-19 & 2026-05-22 & 50 & 55 & 120 \\
Encuesta previa Express & 14 & 13.815 & 22 & 1 & 13.815 & 2019-10-21 & 2026-05-22 & 30 & 30 & 30 \\
Encuesta previa Normal IP & 13 & 23.437 & 22 & 1 & 23.437 & 2019-10-21 & 2026-05-22 & 50 & 60 & 60 \\
Reservas & 29 & 86.562 & 17 & 1 & 86.562 & 2016-08-19 & 2026-05-22 & 20 & 50 & 120 \\
\bottomrule
\end{longtable}
\endgroup


Dentro de cada fuente todos los archivos tienen los mismos encabezados (una sola variante), lo que no demuestra que el cuestionario ni el significado de las etiquetas hayan sido estables en el tiempo. Hay cero filas excedentes por ID repetido, por ID y período repetido, y por duplicado exacto en las cinco fuentes. Las duraciones programadas de las encuestas de satisfacción tienen mediana de 30 minutos (Express) y 55 (Normal); no debe suponerse que toda reserva Normal dura 60 minutos.

\textbf{Estructura de las variables.} Reservas tiene 17 atributos. Las cuatro encuestas repiten esos atributos (el código de servicio se llama «Código servicio» en Reservas y «Código del servicio» en las encuestas), agregan la variable Estado (validez de la encuesta) y añaden preguntas: 4 en las encuestas previas, 15 en la de satisfacción Express y 17 en la de satisfacción Normal.

\textbf{Cobertura de las encuestas previas.} Tienen registros desde el 21 de octubre de 2019 y no cubren todos los períodos de Reservas. Archivos sin registros: Encuesta Express de satisfacción en 201620, 202219, 202319, 202419 y 202519; Encuesta Normal de satisfacción en 202219. Un archivo ausente o vacío no demuestra falta de demanda.

\subsection*{C.2 Filas por período y fuente}

\begingroup\small\setlength{\tabcolsep}{4pt}\renewcommand{\arraystretch}{1.15}
\begin{longtable}{>{\raggedright\arraybackslash}p{\dimexpr 0.131\textwidth-2\tabcolsep\relax}>{\raggedleft\arraybackslash}p{\dimexpr 0.142\textwidth-2\tabcolsep\relax}>{\raggedleft\arraybackslash}p{\dimexpr 0.142\textwidth-2\tabcolsep\relax}>{\raggedleft\arraybackslash}p{\dimexpr 0.142\textwidth-2\tabcolsep\relax}>{\raggedleft\arraybackslash}p{\dimexpr 0.186\textwidth-2\tabcolsep\relax}>{\raggedleft\arraybackslash}p{\dimexpr 0.186\textwidth-2\tabcolsep\relax}}
\toprule
\textbf{Período} & \textbf{Reservas} & \textbf{Encuesta previa Express} & \textbf{Encuesta previa Normal IP} & \textbf{Encuesta de satisfacción Express} & \textbf{Encuesta de satisfacción Normal} \\
\midrule
\endhead
201620 & 2.846 & --- & --- & --- & 1.309 \\
201710 & 3.930 & --- & --- & 781 & 1.320 \\
201719 & 80 & --- & --- & 10 & 38 \\
201720 & 3.372 & --- & --- & 487 & 1.118 \\
201810 & 3.821 & --- & --- & 498 & 1.268 \\
201819 & 132 & --- & --- & 17 & 44 \\
201820 & 4.002 & --- & --- & 594 & 1.251 \\
201910 & 4.495 & --- & --- & 614 & 1.175 \\
201919 & 106 & --- & --- & 21 & 40 \\
201920 & 3.471 & 106 & 157 & 381 & 917 \\
202010 & 4.287 & 1.056 & 1.691 & 22 & 103 \\
202019 & 240 & --- & --- & 6 & 21 \\
202020 & 5.248 & 1.048 & 1.955 & 10 & 31 \\
202110 & 6.077 & 1.157 & 2.274 & 30 & 96 \\
202119 & 110 & --- & --- & 2 & 12 \\
202120 & 4.596 & 1.065 & 1.693 & 76 & 170 \\
202210 & 5.163 & 837 & 2.220 & 35 & 186 \\
202219 & 102 & --- & --- & --- & --- \\
202220 & 5.563 & 1.041 & 2.170 & 26 & 155 \\
202310 & 5.487 & 1.312 & 2.212 & 131 & 277 \\
202319 & 219 & --- & --- & --- & 93 \\
202320 & 5.430 & 1.517 & 2.010 & 780 & 1.520 \\
202410 & 4.990 & 1.340 & 2.228 & 698 & 1.505 \\
202419 & 147 & --- & --- & --- & 73 \\
202420 & 4.870 & 1.359 & 2.302 & 696 & 1.433 \\
202510 & 3.452 & 862 & --- & 461 & 1.203 \\
202519 & 38 & --- & --- & --- & 19 \\
202520 & 2.318 & 570 & 1.402 & 102 & 318 \\
202610 & 1.970 & 545 & 1.123 & 281 & 684 \\
\textbf{Total} & \textbf{86.562} & \textbf{13.815} & \textbf{23.437} & \textbf{6.759} & \textbf{16.379} \\
\bottomrule
\end{longtable}
\endgroup


Los códigos de período tienen seis cifras (año y terminación). Las terminaciones observadas son 10, 19 y 20. El proyecto las interpreta como primer semestre, intersemestral y segundo semestre, interpretación que requiere confirmación del calendario académico; los códigos no se convierten a meses.

\subsection*{C.3 Valores de las variables categóricas de Reservas}

\textbf{Estado de la reserva} (12 valores)

\begingroup\small\setlength{\tabcolsep}{4pt}\renewcommand{\arraystretch}{1.15}
\begin{longtable}{>{\raggedright\arraybackslash}p{\dimexpr 0.538\textwidth-2\tabcolsep\relax}>{\raggedleft\arraybackslash}p{\dimexpr 0.168\textwidth-2\tabcolsep\relax}>{\raggedleft\arraybackslash}p{\dimexpr 0.224\textwidth-2\tabcolsep\relax}}
\toprule
\textbf{Valor} & \textbf{Filas} & \textbf{\% de Reservas} \\
\midrule
\endhead
Finalizada & 38.541 & 44,52\% \\
Cancelada a tiempo & 18.147 & 20,96\% \\
En cola & 10.215 & 11,80\% \\
No asistió & 5.086 & 5,88\% \\
Cola cancelada por reservación & 4.603 & 5,32\% \\
Realizada & 4.105 & 4,74\% \\
Cancelada por sanción & 2.523 & 2,91\% \\
Cancelada & 1.764 & 2,04\% \\
Cancelada con excusa & 784 & 0,91\% \\
En ejecución & 307 & 0,35\% \\
Cancelada y sancionada & 286 & 0,33\% \\
Reservada & 201 & 0,23\% \\
\bottomrule
\end{longtable}
\endgroup


\textbf{Categoría} (11 valores)

\begingroup\small\setlength{\tabcolsep}{4pt}\renewcommand{\arraystretch}{1.15}
\begin{longtable}{>{\raggedright\arraybackslash}p{\dimexpr 0.444\textwidth-2\tabcolsep\relax}>{\raggedleft\arraybackslash}p{\dimexpr 0.209\textwidth-2\tabcolsep\relax}>{\raggedleft\arraybackslash}p{\dimexpr 0.277\textwidth-2\tabcolsep\relax}}
\toprule
\textbf{Valor} & \textbf{Filas} & \textbf{\% de Reservas} \\
\midrule
\endhead
Tutor Presencial & 55.883 & 64,56\% \\
Deprecated, ID 20 & 17.309 & 20,00\% \\
Deprecated, ID 35 & 4.963 & 5,73\% \\
Deprecated, ID 274 & 3.886 & 4,49\% \\
Monitor Remoto & 2.870 & 3,32\% \\
Candidatos & 903 & 1,04\% \\
Monitor Presencial & 427 & 0,49\% \\
Sala & 298 & 0,34\% \\
Deprecated, ID 216 & 15 & 0,02\% \\
Monitor Honores & 7 & 0,01\% \\
Deprecated, ID 276 & 1 & 0,00\% \\
\bottomrule
\end{longtable}
\endgroup


\textbf{Tipo de horario} (10 valores)

\begingroup\small\setlength{\tabcolsep}{4pt}\renewcommand{\arraystretch}{1.15}
\begin{longtable}{>{\raggedright\arraybackslash}p{\dimexpr 0.512\textwidth-2\tabcolsep\relax}>{\raggedleft\arraybackslash}p{\dimexpr 0.179\textwidth-2\tabcolsep\relax}>{\raggedleft\arraybackslash}p{\dimexpr 0.238\textwidth-2\tabcolsep\relax}}
\toprule
\textbf{Valor} & \textbf{Filas} & \textbf{\% de Reservas} \\
\midrule
\endhead
Normal & 49.472 & 57,15\% \\
Express & 25.754 & 29,75\% \\
Normal pico & 8.668 & 10,01\% \\
Entrevistas & 903 & 1,04\% \\
Deprecated, Grupal & 777 & 0,90\% \\
Sala & 298 & 0,34\% \\
Deprecated, Torre Séneca & 268 & 0,31\% \\
Deprecated, Nivelación IP & 256 & 0,30\% \\
Deprecated, Express grupal & 157 & 0,18\% \\
Deprecated, Profesores IP & 9 & 0,01\% \\
\bottomrule
\end{longtable}
\endgroup


\textbf{Servicio} (11 valores)

\begingroup\small\setlength{\tabcolsep}{4pt}\renewcommand{\arraystretch}{1.15}
\begin{longtable}{>{\raggedright\arraybackslash}p{\dimexpr 0.505\textwidth-2\tabcolsep\relax}>{\raggedleft\arraybackslash}p{\dimexpr 0.182\textwidth-2\tabcolsep\relax}>{\raggedleft\arraybackslash}p{\dimexpr 0.243\textwidth-2\tabcolsep\relax}}
\toprule
\textbf{Valor} & \textbf{Filas} & \textbf{\% de Reservas} \\
\midrule
\endhead
IP & 53.272 & 61,54\% \\
Deprecated, APO 1 & 21.136 & 24,42\% \\
EDA & 6.127 & 7,08\% \\
Deprecated, APO 2 & 3.922 & 4,53\% \\
Entrevista & 903 & 1,04\% \\
Deprecated, Apoyo & 361 & 0,42\% \\
Sala & 298 & 0,34\% \\
Deprecated, IP Nivelación & 275 & 0,32\% \\
Deprecated, DPOO & 252 & 0,29\% \\
Deprecated, Profesores & 9 & 0,01\% \\
IP Honores & 7 & 0,01\% \\
\bottomrule
\end{longtable}
\endgroup


\textbf{Código servicio} (11 valores)

\begingroup\small\setlength{\tabcolsep}{4pt}\renewcommand{\arraystretch}{1.15}
\begin{longtable}{>{\raggedright\arraybackslash}p{\dimexpr 0.338\textwidth-2\tabcolsep\relax}>{\raggedleft\arraybackslash}p{\dimexpr 0.254\textwidth-2\tabcolsep\relax}>{\raggedleft\arraybackslash}p{\dimexpr 0.338\textwidth-2\tabcolsep\relax}}
\toprule
\textbf{Valor} & \textbf{Filas} & \textbf{\% de Reservas} \\
\midrule
\endhead
ISIS-1221 & 53.272 & 61,54\% \\
ISIS 1204 & 21.136 & 24,42\% \\
ISIS-1225 & 6.127 & 7,08\% \\
ISIS 1205 & 3.922 & 4,53\% \\
Entrevista & 903 & 1,04\% \\
ISIS F & 361 & 0,42\% \\
Sala & 298 & 0,34\% \\
ISIS 1221N & 275 & 0,32\% \\
ISIS 1226 & 252 & 0,29\% \\
Profesor & 9 & 0,01\% \\
ISIS-1222 & 7 & 0,01\% \\
\bottomrule
\end{longtable}
\endgroup


\textbf{Tipo de usuario} (6 valores)

\begingroup\small\setlength{\tabcolsep}{4pt}\renewcommand{\arraystretch}{1.15}
\begin{longtable}{>{\raggedright\arraybackslash}p{\dimexpr 0.481\textwidth-2\tabcolsep\relax}>{\raggedleft\arraybackslash}p{\dimexpr 0.193\textwidth-2\tabcolsep\relax}>{\raggedleft\arraybackslash}p{\dimexpr 0.256\textwidth-2\tabcolsep\relax}}
\toprule
\textbf{Valor} & \textbf{Filas} & \textbf{\% de Reservas} \\
\midrule
\endhead
Normal & 82.013 & 94,74\% \\
Deprecated, Séneca & 1.967 & 2,27\% \\
Deprecated, Nivelación & 1.758 & 2,03\% \\
Aspirante & 736 & 0,85\% \\
Profesor IP & 81 & 0,09\% \\
Estudiante Honores & 7 & 0,01\% \\
\bottomrule
\end{longtable}
\endgroup


\textbf{Prioritaria} (2 valores)

\begingroup\small\setlength{\tabcolsep}{4pt}\renewcommand{\arraystretch}{1.15}
\begin{longtable}{>{\raggedright\arraybackslash}p{\dimexpr 0.279\textwidth-2\tabcolsep\relax}>{\raggedleft\arraybackslash}p{\dimexpr 0.279\textwidth-2\tabcolsep\relax}>{\raggedleft\arraybackslash}p{\dimexpr 0.371\textwidth-2\tabcolsep\relax}}
\toprule
\textbf{Valor} & \textbf{Filas} & \textbf{\% de Reservas} \\
\midrule
\endhead
no & 80.548 & 93,05\% \\
sí & 6.014 & 6,95\% \\
\bottomrule
\end{longtable}
\endgroup


Las etiquetas que empiezan con «Deprecated» son categorías obsoletas del origen. Se conservan sin unirlas a categorías actuales porque no hay un catálogo histórico que las relacione. Reservas tiene además 213 etiquetas de programa (174 tras normalizar la escritura), 9.943 códigos de usuario distintos, 9.932 correos distintos y 742 nombres de prestador no vacíos. Son cardinalidades de valores exportados, no conteos certificados de personas, monitores o programas.

\textbf{Tipo de horario dentro de las encuestas de satisfacción.} El grupo de archivo no determina la modalidad:

\begingroup\small\setlength{\tabcolsep}{4pt}\renewcommand{\arraystretch}{1.15}
\begin{longtable}{>{\raggedright\arraybackslash}p{\dimexpr 0.419\textwidth-2\tabcolsep\relax}>{\raggedright\arraybackslash}p{\dimexpr 0.333\textwidth-2\tabcolsep\relax}>{\raggedleft\arraybackslash}p{\dimexpr 0.178\textwidth-2\tabcolsep\relax}}
\toprule
\textbf{Archivo} & \textbf{Tipo de horario de la reserva} & \textbf{Encuestas} \\
\midrule
\endhead
Encuesta de satisfacción Express & Express & 6.698 \\
Encuesta de satisfacción Express & Deprecated, Grupal & 61 \\
Encuesta de satisfacción Normal & Normal & 14.831 \\
Encuesta de satisfacción Normal & Deprecated, Torre Séneca & 87 \\
Encuesta de satisfacción Normal & Normal pico & 1.461 \\
\bottomrule
\end{longtable}
\endgroup


Por eso las encuestas se segmentan por la modalidad de su reserva y no por el archivo en que vienen.

\subsection*{C.4 Estado de la reserva y marca de llegada}

\begingroup\small\setlength{\tabcolsep}{4pt}\renewcommand{\arraystretch}{1.15}
\begin{longtable}{>{\raggedright\arraybackslash}p{\dimexpr 0.378\textwidth-2\tabcolsep\relax}>{\raggedleft\arraybackslash}p{\dimexpr 0.118\textwidth-2\tabcolsep\relax}>{\raggedleft\arraybackslash}p{\dimexpr 0.144\textwidth-2\tabcolsep\relax}>{\raggedleft\arraybackslash}p{\dimexpr 0.144\textwidth-2\tabcolsep\relax}>{\raggedleft\arraybackslash}p{\dimexpr 0.144\textwidth-2\tabcolsep\relax}}
\toprule
\textbf{Estado de la reserva} & \textbf{Filas} & \textbf{Con llegada} & \textbf{Sin llegada} & \textbf{\% con llegada} \\
\midrule
\endhead
Cancelada & 1.764 & 0 & 1.764 & 0,00\% \\
Cancelada a tiempo & 18.147 & 4 & 18.143 & 0,02\% \\
Cancelada con excusa & 784 & 16 & 768 & 2,04\% \\
Cancelada por sanción & 2.523 & 8 & 2.515 & 0,32\% \\
Cancelada y sancionada & 286 & 0 & 286 & 0,00\% \\
Cola cancelada por reservación & 4.603 & 4 & 4.599 & 0,09\% \\
En cola & 10.215 & 0 & 10.215 & 0,00\% \\
En ejecución & 307 & 307 & 0 & 100,00\% \\
Finalizada & 38.541 & 38.533 & 8 & 99,98\% \\
No asistió & 5.086 & 9 & 5.077 & 0,18\% \\
Realizada & 4.105 & 4.105 & 0 & 100,00\% \\
Reservada & 201 & 154 & 47 & 76,62\% \\
\bottomrule
\end{longtable}
\endgroup


Casos para revisión (no son errores confirmados ni se descartan): 32 estados de cancelación con marca de llegada, 8 «Finalizada» sin llegada y 9 «No asistió» con llegada, en total 49. Los faltantes de llegada y de prestador dependen del estado de la cita. Una marca de llegada no define por sí sola que la sesión se completó, y la falta de llegada no se convierte en inasistencia de manera automática. Para ninguno de los 12 estados existe una definición oficial de si es terminal ni de si pertenece al denominador de asistencia o cancelación; salvo la agrupación R01, todos conservan su etiqueta original.

\subsection*{C.5 Diccionario de variables}

\textbf{Variables comunes} (están en Reservas y se repiten en las cuatro encuestas, salvo indicación):

\begingroup\small\setlength{\tabcolsep}{4pt}\renewcommand{\arraystretch}{1.15}
\begin{longtable}{>{\raggedright\arraybackslash}p{\dimexpr 0.288\textwidth-2\tabcolsep\relax}>{\raggedright\arraybackslash}p{\dimexpr 0.642\textwidth-2\tabcolsep\relax}}
\toprule
\textbf{Variable} & \textbf{Descripción inicial} \\
\midrule
\endhead
ID & Identificador del evento de reserva. Único en cada grupo; todas las encuestas lo comparten con Reservas. \\
Fecha inicio & Inicio programado del evento, interpretado a partir del encabezado y del formato Excel. \\
Fecha fin & Fin programado del evento. No equivale a una salida real registrada. \\
Fecha de llegada & Marca temporal de llegada según el encabezado; confirmar procedimiento de captura y excepciones. \\
Fecha de devolución & Campo temporal sin valores en todos los archivos; su significado operativo debe confirmarse. \\
Estado de la reserva & Estado del evento de reserva, con 12 etiquetas en la fuente principal. No confundir con Estado de la encuesta. \\
Prioritaria & Indicador sí/no de prioridad de la reserva; confirmar regla de asignación. \\
Categoría & Categoría de recurso o prestación, como Tutor Presencial; incluye identificadores históricos Deprecated. \\
Tipo de horario & Modalidad o tipo de agenda: Normal, Express, Normal pico y otras modalidades históricas. \\
Servicio & Servicio o curso asociado: IP, EDA y etiquetas históricas, entre otros. \\
Código servicio & Código de servicio en Reservas; corresponde al campo Código del servicio en encuestas. \\
Código del servicio & Código de servicio en encuestas; homologar el nombre con Código servicio sin modificar el original. \\
Usuario & Nombre del usuario que reserva; no utilizar como clave única de persona. \\
Correo usuario & Correo del usuario; conservar como atributo, sujeto a validación de identidad. \\
Código usuario & Identificador de usuario exportado como texto. Sus 9.943 valores distintos en Reservas no demuestran 9.943 estudiantes únicos. \\
Tipo de usuario & Clasificación del usuario en Bookeau; no corresponde al programa académico. \\
Programa & Programa reportado para el usuario. Presenta variantes de escritura; no equivale a una lista homologada de programas. \\
Prestador del servicio & Prestador o tutor asociado según el encabezado. Frecuentemente vacío en eventos sin atención; confirmar lógica de asignación. \\
Estado & Estado de validez de la respuesta de encuesta: Válida o Inválida. No es el estado de asistencia. \\
\bottomrule
\end{longtable}
\endgroup


Los significados se infieren de los encabezados y del instructivo de la fuente, salvo los vínculos por ID, que sí se verificaron. Deben confirmarse con el responsable de CupiTaller.

\textbf{Preguntas de la Encuesta de satisfacción Express}

\begingroup\small\setlength{\tabcolsep}{4pt}\renewcommand{\arraystretch}{1.15}
\begin{longtable}{>{\raggedright\arraybackslash}p{\dimexpr 0.426\textwidth-2\tabcolsep\relax}>{\raggedright\arraybackslash}p{\dimexpr 0.184\textwidth-2\tabcolsep\relax}>{\raggedleft\arraybackslash}p{\dimexpr 0.117\textwidth-2\tabcolsep\relax}>{\raggedleft\arraybackslash}p{\dimexpr 0.110\textwidth-2\tabcolsep\relax}>{\raggedright\arraybackslash}p{\dimexpr 0.093\textwidth-2\tabcolsep\relax}}
\toprule
\textbf{Pregunta} & \textbf{Tipo} & \textbf{Respuestas no vacías} & \textbf{\% faltantes} & \textbf{Primer período con valores} \\
\midrule
\endhead
¿El tiempo fue adecuado para resolver su problema o duda puntual? & categoría o respuesta & 6.759 & 0,0\% & 201710 \\
La mayor parte del tiempo de la tutoría lo dediqué a & categoría o respuesta & 3.529 & 47,8\% & 201920 \\
Califique la ayuda que le dio su tutor & categoría o respuesta & 4.983 & 26,3\% & 201819 \\
El tutor me explicó cómo solucionar mi duda/problema/proyecto & categoría o respuesta & 4.983 & 26,3\% & 201819 \\
El tutor aclaró y reforzó conceptos & categoría o respuesta & 4.983 & 26,3\% & 201819 \\
El tutor me brindó herramientas adicionales para trabajar por mi cuenta & categoría o respuesta & 4.983 & 26,3\% & 201819 \\
El tutor fue receptivo frente a mis inquietudes durante la tutoría & categoría o respuesta & 5.448 & 19,4\% & 201810 \\
El tutor es capaz de expresar claramente su conocimiento sobre el tema o temas trabajados & categoría o respuesta & 6.263 & 7,3\% & 201710 \\
Durante la tutoría, el tutor hizo uso de alguna herramienta de inteligencia artificial, tal como pero no limitándose a ChatGPT, Bard/Gemini o Bing AI. & categoría o respuesta & 2.703 & 60,0\% & 202210 \\
Indique cuál fue la herramienta usada durante la tutoría (si sabe cuál es) y el uso que tuvo la misma en la tutoría. & texto libre & 4 & 99,9\% & 202520 \\
¿Cuál considera que es su comprensión del tema actual del curso? & categoría o respuesta & 1.910 & 71,7\% & 202210 \\
¿Considera que el material del curso permite un buen proceso de aprendizaje? & categoría o respuesta & 1.910 & 71,7\% & 202210 \\
Comente aspectos positivos y negativos del tutor y la tutoría. Recuerde que las respuestas van a ser anónimas. & texto libre & 5.453 & 19,3\% & 201720 \\
Escriba en este recuadro cualquier observación o comentario relacionado al curso y sus contenidos. & texto libre & 564 & 91,7\% & 202320 \\
Sugerencias, reclamos y observaciones & texto libre & 1.730 & 74,4\% & 201710 \\
\bottomrule
\end{longtable}
\endgroup


\textbf{Preguntas de la Encuesta de satisfacción Normal}

\begingroup\small\setlength{\tabcolsep}{4pt}\renewcommand{\arraystretch}{1.15}
\begin{longtable}{>{\raggedright\arraybackslash}p{\dimexpr 0.426\textwidth-2\tabcolsep\relax}>{\raggedright\arraybackslash}p{\dimexpr 0.184\textwidth-2\tabcolsep\relax}>{\raggedleft\arraybackslash}p{\dimexpr 0.117\textwidth-2\tabcolsep\relax}>{\raggedleft\arraybackslash}p{\dimexpr 0.110\textwidth-2\tabcolsep\relax}>{\raggedright\arraybackslash}p{\dimexpr 0.093\textwidth-2\tabcolsep\relax}}
\toprule
\textbf{Pregunta} & \textbf{Tipo} & \textbf{Respuestas no vacías} & \textbf{\% faltantes} & \textbf{Primer período con valores} \\
\midrule
\endhead
La tutoría a la que asistí me permitió entender que la solución de un problema usando herramientas de programación es un proceso que implica leer el problema, entender los requerimientos, planear un algoritmo, escribir el código y probar lo implementado & categoría o respuesta & 16.347 & 0,2\% & 201620 \\
La tutoría a la que asistí se enfocó en el desarrollo de mis habilidades de análisis y no solo en la corrección de mi código & categoría o respuesta & 16.359 & 0,1\% & 201620 \\
La mayor parte del tiempo de la tutoría lo dediqué a & categoría o respuesta & 16.366 & 0,1\% & 201620 \\
Aproveché el tiempo de la tutoría haciendo preguntas puntuales al tutor y participando de manera proactiva en el desarrollo de las mismas & categoría o respuesta & 16.364 & 0,1\% & 201620 \\
Califique la ayuda que le dio su tutor & categoría o respuesta & 11.326 & 30,9\% & 201819 \\
El tutor me explicó cómo solucionar mi duda/problema/proyecto & categoría o respuesta & 11.326 & 30,9\% & 201819 \\
El tutor aclaró y reforzó conceptos & categoría o respuesta & 11.326 & 30,9\% & 201819 \\
El tutor me brindó herramientas adicionales para trabajar por mi cuenta & categoría o respuesta & 16.355 & 0,1\% & 201620 \\
El tutor fue receptivo frente a mis inquietudes durante la tutoría & categoría o respuesta & 16.360 & 0,1\% & 201620 \\
El tutor es capaz de expresar claramente su conocimiento sobre el tema o temas trabajados & categoría o respuesta & 16.366 & 0,1\% & 201620 \\
Durante la tutoría, el tutor hizo uso de alguna herramienta de inteligencia artificial, tal como pero no limitándose a ChatGPT, Bard/Gemini o Bing AI. & categoría o respuesta & 6.074 & 62,9\% & 202120 \\
Indique cuál fue la herramienta usada durante la tutoría (si sabe cuál es) y el uso que tuvo la misma en la tutoría. & texto libre & 22 & 99,9\% & 202510 \\
¿Cuál considera que es su comprensión del tema actual del curso? & categoría o respuesta & 4.442 & 72,9\% & 202120 \\
¿Considera que el material del curso permite un buen proceso de aprendizaje? & categoría o respuesta & 4.442 & 72,9\% & 202120 \\
Comente aspectos positivos y negativos del tutor y la tutoría. Recuerde que las respuestas van a ser anónimas. & texto libre & 12.534 & 23,5\% & 201720 \\
Escriba en este recuadro cualquier observación o comentario relacionado al curso y sus contenidos. & texto libre & 1.329 & 91,9\% & 202120 \\
Sugerencias, reclamos y observaciones & texto libre & 4.053 & 75,3\% & 201620 \\
\bottomrule
\end{longtable}
\endgroup


\textbf{Preguntas de la Encuesta previa Express}

\begingroup\small\setlength{\tabcolsep}{4pt}\renewcommand{\arraystretch}{1.15}
\begin{longtable}{>{\raggedright\arraybackslash}p{\dimexpr 0.426\textwidth-2\tabcolsep\relax}>{\raggedright\arraybackslash}p{\dimexpr 0.184\textwidth-2\tabcolsep\relax}>{\raggedleft\arraybackslash}p{\dimexpr 0.117\textwidth-2\tabcolsep\relax}>{\raggedleft\arraybackslash}p{\dimexpr 0.110\textwidth-2\tabcolsep\relax}>{\raggedright\arraybackslash}p{\dimexpr 0.093\textwidth-2\tabcolsep\relax}}
\toprule
\textbf{Pregunta} & \textbf{Tipo} & \textbf{Respuestas no vacías} & \textbf{\% faltantes} & \textbf{Primer período con valores} \\
\midrule
\endhead
La mayoría del tiempo de la tutoría que estoy reservando lo usaré para & categoría o respuesta & 13.815 & 0,0\% & 201920 \\
Entiendo que durante la tutoría no debo usar mi celular & categoría o respuesta & 13.815 & 0,0\% & 201920 \\
Acepto que si quiero trabajar en una tarea en senecode debo tener una solución probada y que pueda explicar, en caso de no cumplir estas condiciones se configurará un comportamiento inadecuado & categoría o respuesta & 13.815 & 0,0\% & 201920 \\
Acepto que he leído los términos de uso de CupiTaller que están disponibles en la URL https://cupitaller.uniandes.edu.co/terminos-de-uso/ & categoría o respuesta & 13.815 & 0,0\% & 201920 \\
\bottomrule
\end{longtable}
\endgroup


\textbf{Preguntas de la Encuesta previa Normal IP}

\begingroup\small\setlength{\tabcolsep}{4pt}\renewcommand{\arraystretch}{1.15}
\begin{longtable}{>{\raggedright\arraybackslash}p{\dimexpr 0.426\textwidth-2\tabcolsep\relax}>{\raggedright\arraybackslash}p{\dimexpr 0.184\textwidth-2\tabcolsep\relax}>{\raggedleft\arraybackslash}p{\dimexpr 0.117\textwidth-2\tabcolsep\relax}>{\raggedleft\arraybackslash}p{\dimexpr 0.110\textwidth-2\tabcolsep\relax}>{\raggedright\arraybackslash}p{\dimexpr 0.093\textwidth-2\tabcolsep\relax}}
\toprule
\textbf{Pregunta} & \textbf{Tipo} & \textbf{Respuestas no vacías} & \textbf{\% faltantes} & \textbf{Primer período con valores} \\
\midrule
\endhead
La mayoría del tiempo de la tutoría que estoy reservando lo usaré para & categoría o respuesta & 23.437 & 0,0\% & 201920 \\
Entiendo que durante la tutoría no debo usar mi celular & categoría o respuesta & 23.437 & 0,0\% & 201920 \\
Acepto que si quiero trabajar en una tarea en senecode debo tener una solución probada y que pueda explicar, en caso de no cumplir estas condiciones se configurará un comportamiento inadecuado & categoría o respuesta & 23.437 & 0,0\% & 201920 \\
Acepto que he leído los términos de uso de CupiTaller que están disponibles en la URL https://cupitaller.uniandes.edu.co/terminos-de-uso/ & categoría o respuesta & 23.437 & 0,0\% & 201920 \\
\bottomrule
\end{longtable}
\endgroup


El primer período con valores no prueba cuándo se creó una pregunta: puede reflejar cambios del formulario, de las etiquetas o de la exportación. Las encuestas previas piden el uso previsto del tiempo de la tutoría y la aceptación de tres condiciones: no usar el celular, tener una solución probada que se pueda explicar si se trabaja una tarea, y haber leído los términos de uso.

\subsection*{C.6 Campos de texto libre}

\begingroup\footnotesize\setlength{\tabcolsep}{4pt}\renewcommand{\arraystretch}{1.15}
\begin{longtable}{>{\raggedright\arraybackslash}p{\dimexpr 0.194\textwidth-2\tabcolsep\relax}>{\raggedright\arraybackslash}p{\dimexpr 0.321\textwidth-2\tabcolsep\relax}>{\raggedleft\arraybackslash}p{\dimexpr 0.064\textwidth-2\tabcolsep\relax}>{\raggedleft\arraybackslash}p{\dimexpr 0.083\textwidth-2\tabcolsep\relax}>{\raggedleft\arraybackslash}p{\dimexpr 0.083\textwidth-2\tabcolsep\relax}>{\raggedleft\arraybackslash}p{\dimexpr 0.058\textwidth-2\tabcolsep\relax}>{\raggedleft\arraybackslash}p{\dimexpr 0.070\textwidth-2\tabcolsep\relax}>{\raggedleft\arraybackslash}p{\dimexpr 0.058\textwidth-2\tabcolsep\relax}}
\toprule
\textbf{Encuesta} & \textbf{Campo} & \textbf{No vacías} & \textbf{Cobertura} & \textbf{Textos distintos} & \textbf{Largo mín.} & \textbf{Largo mediana} & \textbf{Largo máx.} \\
\midrule
\endhead
Encuesta de satisfacción Express & Indique cuál fue la herramienta usada durante la tutoría (si sabe cuál es) y el uso que tuvo la misma en la tutoría. & 4 & 0,06\% & 3 & 6 & 16,0 & 238 \\
Encuesta de satisfacción Express & Comente aspectos positivos y negativos del tutor y la tutoría. Recuerde que las respuestas van a ser anónimas. & 5.453 & 80,68\% & 3.878 & 1 & 26,0 & 550 \\
Encuesta de satisfacción Express & Escriba en este recuadro cualquier observación o comentario relacionado al curso y sus contenidos. & 564 & 8,34\% & 265 & 1 & 7,0 & 339 \\
Encuesta de satisfacción Express & Sugerencias, reclamos y observaciones & 1.730 & 25,60\% & 655 & 1 & 7,0 & 450 \\
Encuesta de satisfacción Normal & Indique cuál fue la herramienta usada durante la tutoría (si sabe cuál es) y el uso que tuvo la misma en la tutoría. & 22 & 0,13\% & 20 & 1 & 14,5 & 90 \\
Encuesta de satisfacción Normal & Comente aspectos positivos y negativos del tutor y la tutoría. Recuerde que las respuestas van a ser anónimas. & 12.534 & 76,52\% & 8.526 & 1 & 29,0 & 1912 \\
Encuesta de satisfacción Normal & Escriba en este recuadro cualquier observación o comentario relacionado al curso y sus contenidos. & 1.329 & 8,11\% & 580 & 1 & 7,0 & 878 \\
Encuesta de satisfacción Normal & Sugerencias, reclamos y observaciones & 4.053 & 24,75\% & 1.411 & 1 & 7,0 & 2281 \\
\bottomrule
\end{longtable}
\endgroup


Hay 17.987 respuestas no vacías al comentario sobre tutor y tutoría entre las dos encuestas de satisfacción (largo mediano de 26 y 29 caracteres). Tener texto no garantiza un comentario sustantivo: los conteos excluyen vacíos y espacios, pero no clasifican respuestas como «no aplica», signos o frases genéricas. Tras excluir las encuestas inválidas (regla R02) quedan 17.829 comentarios con contenido.

\subsection*{C.7 Consistencia entre Reservas y encuestas}

Diferencias entre el valor de la encuesta y el de su reserva (por ID y período), solo en pares no vacíos; se tolera un segundo de diferencia en fechas.

\begingroup\small\setlength{\tabcolsep}{4pt}\renewcommand{\arraystretch}{1.15}
\begin{longtable}{>{\raggedright\arraybackslash}p{\dimexpr 0.266\textwidth-2\tabcolsep\relax}>{\raggedright\arraybackslash}p{\dimexpr 0.141\textwidth-2\tabcolsep\relax}>{\raggedright\arraybackslash}p{\dimexpr 0.153\textwidth-2\tabcolsep\relax}>{\raggedright\arraybackslash}p{\dimexpr 0.185\textwidth-2\tabcolsep\relax}>{\raggedright\arraybackslash}p{\dimexpr 0.185\textwidth-2\tabcolsep\relax}}
\toprule
\textbf{Variable} & \textbf{Encuesta previa Express} & \textbf{Encuesta previa Normal IP} & \textbf{Encuesta de satisfacción Express} & \textbf{Encuesta de satisfacción Normal} \\
\midrule
\endhead
ID & 0 (0,00\%) & 0 (0,00\%) & 0 (0,00\%) & 0 (0,00\%) \\
Fecha inicio & 0 (0,00\%) & 0 (0,00\%) & 0 (0,00\%) & 0 (0,00\%) \\
Fecha fin & 0 (0,00\%) & 0 (0,00\%) & 0 (0,00\%) & 0 (0,00\%) \\
Fecha de llegada & 0 (0,00\%) & 0 (0,00\%) & 0 (0,00\%) & 0 (0,00\%) \\
Estado de la reserva & 15 (0,11\%) & 36 (0,15\%) & 5 (0,07\%) & 17 (0,10\%) \\
Prioritaria & 0 (0,00\%) & 0 (0,00\%) & 0 (0,00\%) & 0 (0,00\%) \\
Categoría & 0 (0,00\%) & 0 (0,00\%) & 0 (0,00\%) & 0 (0,00\%) \\
Tipo de horario & 0 (0,00\%) & 0 (0,00\%) & 0 (0,00\%) & 0 (0,00\%) \\
Servicio & 0 (0,00\%) & 0 (0,00\%) & 0 (0,00\%) & 0 (0,00\%) \\
Código servicio & 0 (0,00\%) & 0 (0,00\%) & 0 (0,00\%) & 0 (0,00\%) \\
Usuario & 8 (0,06\%) & 23 (0,10\%) & 1 (0,01\%) & 3 (0,02\%) \\
Correo usuario & 0 (0,00\%) & 0 (0,00\%) & 0 (0,00\%) & 0 (0,00\%) \\
Código usuario & 15 (0,11\%) & 10 (0,04\%) & 5 (0,07\%) & 10 (0,06\%) \\
Tipo de usuario & 0 (0,00\%) & 0 (0,00\%) & 0 (0,00\%) & 0 (0,00\%) \\
Programa & 30 (0,22\%) & 20 (0,09\%) & 12 (0,18\%) & 35 (0,21\%) \\
Prestador del servicio & 82 (1,07\%) & 144 (0,96\%) & 67 (0,99\%) & 257 (1,57\%) \\
\bottomrule
\end{longtable}
\endgroup


Las fechas, la categoría, el tipo de horario, el servicio y el código de servicio no tienen diferencias. Las diferencias en estado, usuario, programa y prestador no demuestran errores por sí solas: podrían deberse a actualizaciones posteriores o a variantes de representación. Falta definir con el responsable una fuente de referencia y una regla de conciliación antes de sobrescribir atributos.

\textbf{Pares de estados que difieren entre la reserva y la encuesta}

\begingroup\small\setlength{\tabcolsep}{4pt}\renewcommand{\arraystretch}{1.15}
\begin{longtable}{>{\raggedright\arraybackslash}p{\dimexpr 0.433\textwidth-2\tabcolsep\relax}>{\raggedright\arraybackslash}p{\dimexpr 0.198\textwidth-2\tabcolsep\relax}>{\raggedright\arraybackslash}p{\dimexpr 0.171\textwidth-2\tabcolsep\relax}>{\raggedleft\arraybackslash}p{\dimexpr 0.128\textwidth-2\tabcolsep\relax}}
\toprule
\textbf{Encuesta} & \textbf{Estado en Reservas} & \textbf{Estado en la encuesta} & \textbf{Filas} \\
\midrule
\endhead
Encuesta de satisfacción Express & En ejecución & Realizada & 1 \\
Encuesta de satisfacción Express & Finalizada & Realizada & 1 \\
Encuesta de satisfacción Express & Realizada & Finalizada & 3 \\
Encuesta de satisfacción Normal & En ejecución & Realizada & 1 \\
Encuesta de satisfacción Normal & Finalizada & Realizada & 9 \\
Encuesta de satisfacción Normal & Realizada & Finalizada & 7 \\
Encuesta previa Express & En ejecución & Realizada & 5 \\
Encuesta previa Express & Realizada & Finalizada & 10 \\
Encuesta previa Normal IP & En ejecución & Finalizada & 2 \\
Encuesta previa Normal IP & En ejecución & Realizada & 10 \\
Encuesta previa Normal IP & Realizada & Finalizada & 24 \\
\bottomrule
\end{longtable}
\endgroup


\subsection*{C.8 Reglas de fechas y de respuestas}

\begingroup\small\setlength{\tabcolsep}{4pt}\renewcommand{\arraystretch}{1.15}
\begin{longtable}{>{\raggedright\arraybackslash}p{\dimexpr 0.219\textwidth-2\tabcolsep\relax}>{\raggedright\arraybackslash}p{\dimexpr 0.143\textwidth-2\tabcolsep\relax}>{\raggedright\arraybackslash}p{\dimexpr 0.143\textwidth-2\tabcolsep\relax}>{\raggedright\arraybackslash}p{\dimexpr 0.143\textwidth-2\tabcolsep\relax}>{\raggedright\arraybackslash}p{\dimexpr 0.139\textwidth-2\tabcolsep\relax}>{\raggedright\arraybackslash}p{\dimexpr 0.143\textwidth-2\tabcolsep\relax}}
\toprule
\textbf{Regla} & \textbf{Reservas} & \textbf{Encuesta previa Express} & \textbf{Encuesta previa Normal IP} & \textbf{Encuesta de satisfacción Express} & \textbf{Encuesta de satisfacción Normal} \\
\midrule
\endhead
Fecha fin anterior al inicio & 86.562 evaluables / 0 incumplen & 13.815 evaluables / 0 incumplen & 23.437 evaluables / 0 incumplen & 6.759 evaluables / 0 incumplen & 16.379 evaluables / 0 incumplen \\
Devolución anterior a la llegada & 0 evaluables / 0 incumplen & 0 evaluables / 0 incumplen & 0 evaluables / 0 incumplen & 0 evaluables / 0 incumplen & 0 evaluables / 0 incumplen \\
Año del inicio distinto del año del código de período & 86.562 evaluables / 0 incumplen & 13.815 evaluables / 0 incumplen & 23.437 evaluables / 0 incumplen & 6.759 evaluables / 0 incumplen & 16.379 evaluables / 0 incumplen \\
\bottomrule
\end{longtable}
\endgroup


La coherencia entre llegada y devolución no es evaluable porque la fecha de devolución está vacía en todas las filas de las cinco fuentes. Estas comprobaciones son parciales: no validan todas las reglas del negocio.

\begingroup\small\setlength{\tabcolsep}{4pt}\renewcommand{\arraystretch}{1.15}
\begin{longtable}{>{\raggedright\arraybackslash}p{\dimexpr 0.191\textwidth-2\tabcolsep\relax}>{\raggedright\arraybackslash}p{\dimexpr 0.219\textwidth-2\tabcolsep\relax}>{\raggedright\arraybackslash}p{\dimexpr 0.290\textwidth-2\tabcolsep\relax}>{\raggedleft\arraybackslash}p{\dimexpr 0.081\textwidth-2\tabcolsep\relax}>{\raggedleft\arraybackslash}p{\dimexpr 0.069\textwidth-2\tabcolsep\relax}>{\raggedleft\arraybackslash}p{\dimexpr 0.081\textwidth-2\tabcolsep\relax}}
\toprule
\textbf{Encuesta} & \textbf{Variable} & \textbf{Regla} & \textbf{No vacíos evaluados} & \textbf{Fuera de dominio} & \textbf{Faltantes} \\
\midrule
\endhead
Encuesta de satisfacción Express & Califique la ayuda que le dio su tutor & valor entero entre 1 y 5 (dominio candidato observado) & 4.983 & 0 & 1.776 \\
Encuesta de satisfacción Normal & Califique la ayuda que le dio su tutor & valor entero entre 1 y 5 (dominio candidato observado) & 11.326 & 0 & 5.053 \\
\bottomrule
\end{longtable}
\endgroup


El dominio 1 a 5 de la calificación se deriva de lo observado; hay que confirmar el cuestionario y sus anclajes antes de interpretar promedios.

\subsection*{C.9 Estados de reserva y encuestas}

La coordinación indicó que «Finalizada» y «Realizada» se tratan igual y sugirió que la diferencia podría ser una encuesta incompleta. Se comprobó cruzando por ID los estados de Reservas con las encuestas.

\begingroup\small\setlength{\tabcolsep}{4pt}\renewcommand{\arraystretch}{1.15}
\begin{longtable}{>{\raggedright\arraybackslash}p{\dimexpr 0.306\textwidth-2\tabcolsep\relax}>{\raggedleft\arraybackslash}p{\dimexpr 0.117\textwidth-2\tabcolsep\relax}>{\raggedleft\arraybackslash}p{\dimexpr 0.137\textwidth-2\tabcolsep\relax}>{\raggedleft\arraybackslash}p{\dimexpr 0.137\textwidth-2\tabcolsep\relax}>{\raggedleft\arraybackslash}p{\dimexpr 0.127\textwidth-2\tabcolsep\relax}>{\raggedleft\arraybackslash}p{\dimexpr 0.106\textwidth-2\tabcolsep\relax}}
\toprule
\textbf{Estado en Reservas} & \textbf{Eventos} & \textbf{Con encuesta posterior} & \textbf{\% posterior} & \textbf{Con encuesta previa} & \textbf{\% previa} \\
\midrule
\endhead
Cancelada & 1.764 & 0 & 0,00\% & 0 & 0,00\% \\
Cancelada a tiempo & 18.147 & 0 & 0,00\% & 11.496 & 63,35\% \\
Cancelada con excusa & 784 & 5 & 0,64\% & 372 & 47,45\% \\
Cancelada por sanción & 2.523 & 0 & 0,00\% & 277 & 10,98\% \\
Cancelada y sancionada & 286 & 0 & 0,00\% & 0 & 0,00\% \\
Cola cancelada por reservación & 4.603 & 0 & 0,00\% & 0 & 0,00\% \\
En cola & 10.215 & 0 & 0,00\% & 0 & 0,00\% \\
En ejecución & 307 & 2 & 0,65\% & 190 & 61,89\% \\
Finalizada & 38.541 & 21.507 & 55,80\% & 19.427 & 50,41\% \\
No asistió & 5.086 & 0 & 0,00\% & 2.540 & 49,94\% \\
Realizada & 4.105 & 1.620 & 39,46\% & 2.897 & 70,57\% \\
Reservada & 201 & 4 & 1,99\% & 53 & 26,37\% \\
\bottomrule
\end{longtable}
\endgroup


Los porcentajes usan como denominador todos los eventos de cada estado, con todas las modalidades y períodos. Son cobertura en los archivos compartidos, no tasas de respuesta de personas invitadas: una encuesta no exportada podría no haber sido requerida, existir fuera del paquete o no haber sido respondida.

\textbf{Validez y presencia de respuestas en las encuestas de satisfacción, por estado de la reserva}

\begingroup\footnotesize\setlength{\tabcolsep}{4pt}\renewcommand{\arraystretch}{1.15}
\begin{longtable}{>{\raggedright\arraybackslash}p{\dimexpr 0.218\textwidth-2\tabcolsep\relax}>{\raggedright\arraybackslash}p{\dimexpr 0.150\textwidth-2\tabcolsep\relax}>{\raggedleft\arraybackslash}p{\dimexpr 0.093\textwidth-2\tabcolsep\relax}>{\raggedleft\arraybackslash}p{\dimexpr 0.079\textwidth-2\tabcolsep\relax}>{\raggedleft\arraybackslash}p{\dimexpr 0.093\textwidth-2\tabcolsep\relax}>{\raggedleft\arraybackslash}p{\dimexpr 0.093\textwidth-2\tabcolsep\relax}>{\raggedleft\arraybackslash}p{\dimexpr 0.093\textwidth-2\tabcolsep\relax}>{\raggedleft\arraybackslash}p{\dimexpr 0.112\textwidth-2\tabcolsep\relax}}
\toprule
\textbf{Encuesta} & \textbf{Estado en Reservas} & \textbf{Registros} & \textbf{Válidas} & \textbf{Inválidas} & \textbf{Sin ninguna respuesta} & \textbf{Con alguna respuesta} & \textbf{Con calificación} \\
\midrule
\endhead
Encuesta de satisfacción Express & En ejecución & 1 & 1 & 0 & 0 & 1 & 1 \\
Encuesta de satisfacción Express & Finalizada & 6.299 & 6.248 & 51 & 0 & 6.299 & 4.525 \\
Encuesta de satisfacción Express & Realizada & 457 & 455 & 2 & 0 & 457 & 455 \\
Encuesta de satisfacción Express & Reservada & 2 & 2 & 0 & 0 & 2 & 2 \\
Encuesta de satisfacción Normal & Cancelada con excusa & 5 & 2 & 3 & 1 & 4 & 3 \\
Encuesta de satisfacción Normal & En ejecución & 1 & 1 & 0 & 0 & 1 & 1 \\
Encuesta de satisfacción Normal & Finalizada & 15.208 & 15.101 & 107 & 12 & 15.196 & 10.164 \\
Encuesta de satisfacción Normal & Realizada & 1.163 & 1.162 & 1 & 0 & 1.163 & 1.156 \\
Encuesta de satisfacción Normal & Reservada & 2 & 2 & 0 & 0 & 2 & 2 \\
\bottomrule
\end{longtable}
\endgroup


\textbf{Conclusión de la comprobación.} La extracción no confirma que «Finalizada» o «Realizada» identifique si una encuesta fue completada: hay reservas con y sin encuesta en ambos estados, y las etiquetas de invalidez aparecen en ambos. En la encuesta Normal hay 12 registros de «Finalizada» sin ninguna respuesta no vacía; se revisan como casos de cobertura y no se descartan. Tampoco hay una fuente separada de encuestas de tutores, así que esa parte de la explicación no se puede comprobar. Dentro de período, modalidad y servicio se encontraron 103 estratos con ambos estados: en 55 la cobertura de encuesta posterior es mayor en «Finalizada», en 3 mayor en «Realizada» y en 45 igual; entre las modalidades Express y Normal hay 78 estratos con ambos estados. Estos conteos no son una prueba estadística. «Todas las preguntas no vacías» tampoco es un criterio de encuesta completa, porque hay preguntas opcionales, condicionales e históricas cuya obligatoriedad se desconoce.

\subsection*{C.10 Comparación Express y Normal por período y servicio}

Hay 37 combinaciones de período y servicio con calificaciones en las dos modalidades exactas, después de excluir las encuestas inválidas. La coincidencia de pregunta, período y servicio es necesaria para comparar, pero no demuestra que los formularios fueran equivalentes ni elimina sesgos.

\begingroup\footnotesize\setlength{\tabcolsep}{4pt}\renewcommand{\arraystretch}{1.15}
\begin{longtable}{>{\raggedright\arraybackslash}p{\dimexpr 0.095\textwidth-2\tabcolsep\relax}>{\raggedright\arraybackslash}p{\dimexpr 0.216\textwidth-2\tabcolsep\relax}>{\raggedleft\arraybackslash}p{\dimexpr 0.095\textwidth-2\tabcolsep\relax}>{\raggedleft\arraybackslash}p{\dimexpr 0.087\textwidth-2\tabcolsep\relax}>{\raggedleft\arraybackslash}p{\dimexpr 0.095\textwidth-2\tabcolsep\relax}>{\raggedleft\arraybackslash}p{\dimexpr 0.087\textwidth-2\tabcolsep\relax}>{\raggedleft\arraybackslash}p{\dimexpr 0.128\textwidth-2\tabcolsep\relax}>{\raggedleft\arraybackslash}p{\dimexpr 0.128\textwidth-2\tabcolsep\relax}}
\toprule
\textbf{Período} & \textbf{Servicio} & \textbf{n Express} & \textbf{n Normal} & \textbf{Media Express} & \textbf{Media Normal} & \textbf{Completitud Express} & \textbf{Completitud Normal} \\
\midrule
\endhead
201819 & Deprecated, APO 1 & 10 & 31 & 4,80 & 4,94 & 100,0\% & 100,0\% \\
201819 & Deprecated, APO 2 & 7 & 13 & 5,00 & 4,92 & 100,0\% & 100,0\% \\
201820 & Deprecated, APO 1 & 490 & 1086 & 4,86 & 4,88 & 100,0\% & 100,0\% \\
201820 & Deprecated, APO 2 & 50 & 108 & 4,90 & 4,81 & 100,0\% & 100,0\% \\
201820 & Deprecated, Apoyo & 14 & 22 & 4,93 & 4,95 & 100,0\% & 100,0\% \\
201820 & Deprecated, IP Nivelación & 1 & 1 & 5,00 & 5,00 & 100,0\% & 100,0\% \\
201910 & Deprecated, APO 1 & 507 & 1022 & 4,81 & 4,85 & 100,0\% & 100,0\% \\
201910 & Deprecated, APO 2 & 48 & 94 & 4,77 & 4,80 & 100,0\% & 100,0\% \\
201910 & Deprecated, Apoyo & 10 & 3 & 5,00 & 4,67 & 100,0\% & 100,0\% \\
201910 & Deprecated, IP Nivelación & 11 & 14 & 5,00 & 4,79 & 100,0\% & 100,0\% \\
201919 & Deprecated, APO 1 & 16 & 31 & 4,81 & 4,90 & 100,0\% & 100,0\% \\
201919 & Deprecated, APO 2 & 5 & 9 & 5,00 & 4,67 & 100,0\% & 100,0\% \\
201920 & Deprecated, APO 1 & 329 & 698 & 4,88 & 4,86 & 100,0\% & 100,0\% \\
201920 & Deprecated, APO 2 & 49 & 130 & 4,92 & 4,86 & 100,0\% & 100,0\% \\
202010 & Deprecated, APO 2 & 22 & 100 & 4,86 & 4,64 & 100,0\% & 100,0\% \\
202019 & Deprecated, APO 2 & 6 & 21 & 4,83 & 4,95 & 100,0\% & 100,0\% \\
202020 & Deprecated, APO 2 & 10 & 31 & 4,70 & 4,90 & 100,0\% & 100,0\% \\
202110 & EDA & 30 & 96 & 4,37 & 4,76 & 100,0\% & 100,0\% \\
202119 & EDA & 2 & 12 & 5,00 & 5,00 & 100,0\% & 100,0\% \\
202120 & EDA & 76 & 169 & 4,64 & 4,68 & 100,0\% & 100,0\% \\
202210 & EDA & 34 & 186 & 4,53 & 4,40 & 100,0\% & 100,0\% \\
202220 & EDA & 24 & 148 & 4,67 & 4,74 & 100,0\% & 100,0\% \\
202220 & IP & 1 & 7 & 5,00 & 4,86 & 100,0\% & 100,0\% \\
202310 & EDA & 124 & 250 & 4,65 & 4,78 & 100,0\% & 100,0\% \\
202310 & IP & 7 & 18 & 5,00 & 4,89 & 100,0\% & 100,0\% \\
202320 & EDA & 138 & 300 & 4,67 & 4,70 & 100,0\% & 100,0\% \\
202320 & IP & 632 & 823 & 4,62 & 4,76 & 100,0\% & 100,0\% \\
202410 & EDA & 37 & 202 & 4,81 & 4,71 & 100,0\% & 100,0\% \\
202410 & IP & 659 & 822 & 4,62 & 4,71 & 100,0\% & 100,0\% \\
202420 & EDA & 12 & 83 & 4,50 & 4,55 & 100,0\% & 100,0\% \\
202420 & IP & 675 & 1002 & 4,60 & 4,75 & 100,0\% & 100,0\% \\
202510 & EDA & 31 & 75 & 4,77 & 4,84 & 100,0\% & 100,0\% \\
202510 & IP & 430 & 923 & 4,61 & 4,70 & 100,0\% & 100,0\% \\
202520 & EDA & 4 & 17 & 5,00 & 4,47 & 100,0\% & 100,0\% \\
202520 & IP & 95 & 267 & 4,61 & 4,79 & 100,0\% & 100,0\% \\
202610 & EDA & 14 & 28 & 4,79 & 4,82 & 100,0\% & 100,0\% \\
202610 & IP & 267 & 627 & 4,69 & 4,73 & 100,0\% & 100,0\% \\
\bottomrule
\end{longtable}
\endgroup


\subsection*{C.11 Capa Silver: reglas, controles y resultados de la limpieza}

Se generaron cinco archivos CSV con las 146.952 filas de origen, sin eliminar ninguna.

\textbf{Reglas aplicadas}

\begingroup\small\setlength{\tabcolsep}{4pt}\renewcommand{\arraystretch}{1.15}
\begin{longtable}{>{\raggedright\arraybackslash}p{\dimexpr 0.265\textwidth-2\tabcolsep\relax}>{\raggedright\arraybackslash}p{\dimexpr 0.665\textwidth-2\tabcolsep\relax}}
\toprule
\textbf{Regla} & \textbf{Resultado} \\
\midrule
\endhead
Nombres de columnas & Nombres comunes para los atributos, incluido \texttt{codigo\_\allowbreak{}servicio} para ambos encabezados de origen. Las preguntas reciben nombres estables derivados de su texto completo. \\
Nulos & Las celdas ausentes, vacías o con solo espacios se guardan vacías. «No aplica» y otras respuestas de texto se mantienen como valores. \\
Fechas & Se leen con los estilos y el sistema de fechas de Excel y se guardan en ISO 8601 con microsegundos, sin inventar zona horaria. \\
Identificadores & El ID de reserva y los códigos se guardan como texto. No se corrigen ni se fusionan identidades. \\
Programas & Columna adicional \texttt{programa\_\allowbreak{}normalizado}: mayúsculas, sin tildes y con espacios compactados (incluidos los no separables). \texttt{programa} conserva el valor original. No se fusionan carreras por significado. \\
Prioridad & \texttt{es\_\allowbreak{}prioritaria} toma true o false a partir de sí o no; \texttt{prioritaria\_\allowbreak{}original} conserva el valor exportado. \\
Calificación & Campo numérico adicional para valores enteros de 1 a 5; el original se conserva; no se imputan las respuestas ausentes. \\
Trazabilidad & Cada fila conserva archivo, hoja, fila física de Excel y período de origen. \\
Calidad & Banderas por fila para casos que requieren revisión; las encuestas conservan sus propios atributos aunque difieran de Reservas. \\
\bottomrule
\end{longtable}
\endgroup


El cambio de formato de una fecha no crea información nueva. El programa normalizado solo homologa la escritura; no es un catálogo académico certificado. Los textos libres, correos, nombres, estados originales y categorías históricas no se reescriben.

\textbf{Controles por fuente}

\begingroup\footnotesize\setlength{\tabcolsep}{4pt}\renewcommand{\arraystretch}{1.15}
\begin{longtable}{>{\raggedright\arraybackslash}p{\dimexpr 0.183\textwidth-2\tabcolsep\relax}>{\raggedleft\arraybackslash}p{\dimexpr 0.060\textwidth-2\tabcolsep\relax}>{\raggedleft\arraybackslash}p{\dimexpr 0.060\textwidth-2\tabcolsep\relax}>{\raggedleft\arraybackslash}p{\dimexpr 0.083\textwidth-2\tabcolsep\relax}>{\raggedleft\arraybackslash}p{\dimexpr 0.078\textwidth-2\tabcolsep\relax}>{\raggedleft\arraybackslash}p{\dimexpr 0.078\textwidth-2\tabcolsep\relax}>{\raggedleft\arraybackslash}p{\dimexpr 0.066\textwidth-2\tabcolsep\relax}>{\raggedleft\arraybackslash}p{\dimexpr 0.078\textwidth-2\tabcolsep\relax}>{\raggedleft\arraybackslash}p{\dimexpr 0.078\textwidth-2\tabcolsep\relax}>{\raggedleft\arraybackslash}p{\dimexpr 0.072\textwidth-2\tabcolsep\relax}>{\raggedright\arraybackslash}p{\dimexpr 0.094\textwidth-2\tabcolsep\relax}}
\toprule
\textbf{Fuente} & \textbf{Filas Bronze} & \textbf{Filas Silver} & \textbf{Eliminadas} & \textbf{ID distintos} & \textbf{Encuestas sin reserva} & \textbf{Con alguna bandera} & \textbf{Incluidas R02} & \textbf{Excluidas R02} & \textbf{Atendida (R01)} & \textbf{Conciliación} \\
\midrule
\endhead
Encuesta de satisfacción Express & 6.759 & 6.759 & 0 & 6.759 & 0 & 143 & 6.706 & 53 & 6.757 & OK \\
Encuesta de satisfacción Normal & 16.379 & 16.379 & 0 & 16.379 & 0 & 427 & 16.268 & 111 & 16.372 & OK \\
Encuesta previa Express & 13.815 & 13.815 & 0 & 13.815 & 0 & 144 & 13.815 & 0 & 7.559 & OK \\
Encuesta previa Normal IP & 23.437 & 23.437 & 0 & 23.437 & 0 & 220 & 23.435 & 2 & 14.782 & OK \\
Reservas & 86.562 & 86.562 & 0 & 86.562 & 0 & 49 & --- & --- & 42.646 & OK \\
\bottomrule
\end{longtable}
\endgroup


\textbf{Programas antes y después de normalizar la escritura}

\begingroup\small\setlength{\tabcolsep}{4pt}\renewcommand{\arraystretch}{1.15}
\begin{longtable}{>{\raggedright\arraybackslash}p{\dimexpr 0.471\textwidth-2\tabcolsep\relax}>{\raggedleft\arraybackslash}p{\dimexpr 0.215\textwidth-2\tabcolsep\relax}>{\raggedleft\arraybackslash}p{\dimexpr 0.243\textwidth-2\tabcolsep\relax}}
\toprule
\textbf{Fuente} & \textbf{Etiquetas originales} & \textbf{Etiquetas normalizadas} \\
\midrule
\endhead
Encuesta de satisfacción Express & 131 & 109 \\
Encuesta de satisfacción Normal & 140 & 119 \\
Encuesta previa Express & 143 & 115 \\
Encuesta previa Normal IP & 122 & 91 \\
Reservas & 213 & 174 \\
\bottomrule
\end{longtable}
\endgroup


La reducción de etiquetas no representa una reducción de personas ni certifica equivalencias entre titulaciones. Ejemplo: «INGENIERIA ELECTRONICA», «Ingeniería Electrónica» e «INGENIERÍA ELECTRONICA» pasan a la misma etiqueta normalizada.

\textbf{Banderas de calidad}

\begingroup\small\setlength{\tabcolsep}{4pt}\renewcommand{\arraystretch}{1.15}
\begin{longtable}{>{\raggedright\arraybackslash}p{\dimexpr 0.387\textwidth-2\tabcolsep\relax}>{\raggedright\arraybackslash}p{\dimexpr 0.378\textwidth-2\tabcolsep\relax}>{\raggedleft\arraybackslash}p{\dimexpr 0.165\textwidth-2\tabcolsep\relax}}
\toprule
\textbf{Fuente} & \textbf{Bandera} & \textbf{Registros} \\
\midrule
\endhead
Encuesta de satisfacción Express & \texttt{atributos\_\allowbreak{}difieren\_\allowbreak{}de\_\allowbreak{}reserva} & 88 \\
Encuesta de satisfacción Express & \texttt{encuesta\_\allowbreak{}marcada\_\allowbreak{}invalida\_\allowbreak{}en\_\allowbreak{}origen} & 53 \\
Encuesta de satisfacción Express & \texttt{revisar\_\allowbreak{}finalizada\_\allowbreak{}sin\_\allowbreak{}llegada} & 2 \\
Encuesta de satisfacción Normal & \texttt{atributos\_\allowbreak{}difieren\_\allowbreak{}de\_\allowbreak{}reserva} & 311 \\
Encuesta de satisfacción Normal & \texttt{encuesta\_\allowbreak{}marcada\_\allowbreak{}invalida\_\allowbreak{}en\_\allowbreak{}origen} & 111 \\
Encuesta de satisfacción Normal & \texttt{revisar\_\allowbreak{}cancelacion\_\allowbreak{}con\_\allowbreak{}llegada} & 2 \\
Encuesta de satisfacción Normal & \texttt{revisar\_\allowbreak{}finalizada\_\allowbreak{}sin\_\allowbreak{}llegada} & 4 \\
Encuesta previa Express & \texttt{atributos\_\allowbreak{}difieren\_\allowbreak{}de\_\allowbreak{}reserva} & 144 \\
Encuesta previa Normal IP & \texttt{atributos\_\allowbreak{}difieren\_\allowbreak{}de\_\allowbreak{}reserva} & 218 \\
Encuesta previa Normal IP & \texttt{encuesta\_\allowbreak{}marcada\_\allowbreak{}invalida\_\allowbreak{}en\_\allowbreak{}origen} & 2 \\
Reservas & \texttt{revisar\_\allowbreak{}cancelacion\_\allowbreak{}con\_\allowbreak{}llegada} & 32 \\
Reservas & \texttt{revisar\_\allowbreak{}finalizada\_\allowbreak{}sin\_\allowbreak{}llegada} & 8 \\
Reservas & \texttt{revisar\_\allowbreak{}no\_\allowbreak{}asistio\_\allowbreak{}con\_\allowbreak{}llegada} & 9 \\
\bottomrule
\end{longtable}
\endgroup


Un registro puede activar varias banderas, así que las frecuencias no se suman como registros únicos. Las 166 encuestas inválidas se conservan en Silver con \texttt{incluir\_\allowbreak{}encuesta\_\allowbreak{}en\_\allowbreak{}analisis = false}. La ausencia total de fecha de devolución no se imputa con la fecha fin, y los faltantes de llegada no se convierten en «No asistió». El resultado «OK» de la conciliación estructural confirma filas y vínculos; no significa que todos los estados o valores sean correctos para el negocio.

\textbf{Campos que agrega Silver}

\begingroup\small\setlength{\tabcolsep}{4pt}\renewcommand{\arraystretch}{1.15}
\begin{longtable}{>{\raggedright\arraybackslash}p{\dimexpr 0.248\textwidth-2\tabcolsep\relax}>{\raggedright\arraybackslash}p{\dimexpr 0.182\textwidth-2\tabcolsep\relax}>{\raggedright\arraybackslash}p{\dimexpr 0.500\textwidth-2\tabcolsep\relax}}
\toprule
\textbf{Campo} & \textbf{Tipo} & \textbf{Significado} \\
\midrule
\endhead
\texttt{periodo\_\allowbreak{}origen} & texto & Código de período derivado del nombre del archivo; no se convierte en mes o fecha. \\
\texttt{archivo\_\allowbreak{}origen} & texto & Ruta relativa del Excel original. \\
\texttt{hoja\_\allowbreak{}origen} & texto & Nombre de la hoja del libro. \\
\texttt{fila\_\allowbreak{}excel} & entero & Número de fila física Excel; las celdas ausentes no desplazan columnas. \\
\texttt{programa\_\allowbreak{}normalizado} & texto & Etiqueta adicional en mayúsculas, sin tildes y con espacios compactados; programa conserva el original no vacío. \\
\texttt{es\_\allowbreak{}prioritaria} & booleano nullable & true/false a partir de sí/no; vacío si el indicador no se reconoce. \\
\texttt{duracion\_\allowbreak{}programada\_\allowbreak{}minutos} & decimal nullable & Diferencia no negativa entre fecha\_fin y fecha\_inicio; no representa duración efectiva. \\
\texttt{banderas\_\allowbreak{}calidad\_\allowbreak{}json} & array JSON & Motivos para revisar; [] indica que no se activó una regla de esta etapa, no calidad garantizada. \\
\texttt{marcada\_\allowbreak{}valida\_\allowbreak{}en\_\allowbreak{}origen} & booleano nullable & Solo en encuestas; copia de la etiqueta Válida/Inválida como indicador, sin redefinir su significado. \\
\texttt{reserva\_\allowbreak{}encontrada} & booleano & Solo en encuestas; existencia de id\_reserva en la fuente Reservas. \\
\texttt{diferencias\_\allowbreak{}con\_\allowbreak{}reserva\_\allowbreak{}json} & array JSON & Solo en encuestas; campos no vacíos que difieren de la reserva asociada, tolerando un segundo en fechas. \\
\texttt{calificacion\_\allowbreak{}ayuda\_\allowbreak{}tutor} & entero nullable & Solo en satisfacción; copia numérica utilizable del valor observado 1--5; original conservado en calificacion\_ayuda\_tutor\_original. \\
\texttt{calificacion\_\allowbreak{}en\_\allowbreak{}dominio} & booleano nullable & Solo en satisfacción; true si el valor es 1--5, false si es otro valor no vacío, vacío si no hay respuesta. \\
\texttt{estado\_\allowbreak{}reserva\_\allowbreak{}analitico} & texto & R01: Finalizada y Realizada $\rightarrow$ Atendida, según regla de trabajo aportada por María Alejandra Pérez. Otros estados mantienen su etiqueta; estado\_reserva conserva siempre el original. No determina encuesta completada. \\
\texttt{incluir\_\allowbreak{}encuesta\_\allowbreak{}en\_\allowbreak{}analisis} & booleano nullable & R02: true para Válida, false para Inválida; vacío si estado desconocido. Aplicar a métricas y análisis de respuestas, no a conteos de auditoría. \\
\texttt{motivo\_\allowbreak{}exclusion\_\allowbreak{}encuesta} & texto nullable & Motivo explícito de exclusión según R02 o estado desconocido; vacío cuando se incluye la encuesta. \\
\bottomrule
\end{longtable}
\endgroup


\textbf{Celdas modificadas por regla (suma de las cinco fuentes)}

\begingroup\small\setlength{\tabcolsep}{4pt}\renewcommand{\arraystretch}{1.15}
\begin{longtable}{>{\raggedright\arraybackslash}p{\dimexpr 0.699\textwidth-2\tabcolsep\relax}>{\raggedleft\arraybackslash}p{\dimexpr 0.231\textwidth-2\tabcolsep\relax}}
\toprule
\textbf{Regla} & \textbf{Celdas afectadas} \\
\midrule
\endhead
Celda vacía o con espacios convertida a nulo & 404.949 \\
Fecha de Excel convertida a ISO 8601 & 382.738 \\
R01: Finalizada y Realizada a Atendida & 88.116 \\
Programa normalizado en columna adicional & 13.393 \\
R02: encuesta inválida excluida del análisis & 166 \\
\bottomrule
\end{longtable}
\endgroup


\subsection*{C.12 Capa Gold: tablas, controles y conciliación}

\begingroup\small\setlength{\tabcolsep}{4pt}\renewcommand{\arraystretch}{1.15}
\begin{longtable}{>{\raggedright\arraybackslash}p{\dimexpr 0.203\textwidth-2\tabcolsep\relax}>{\raggedleft\arraybackslash}p{\dimexpr 0.110\textwidth-2\tabcolsep\relax}>{\raggedright\arraybackslash}p{\dimexpr 0.616\textwidth-2\tabcolsep\relax}}
\toprule
\textbf{Tabla} & \textbf{Filas} & \textbf{Contenido} \\
\midrule
\endhead
\texttt{dim\_\allowbreak{}fecha} & 3.565 & Calendario diario continuo del rango observado (inicio, fin y llegada), con el período académico como atributo \\
\texttt{dim\_\allowbreak{}hora} & 1.440 & Un miembro por minuto del día \\
\texttt{dim\_\allowbreak{}servicio} & 11 & Código y etiqueta original del servicio \\
\texttt{dim\_\allowbreak{}modalidad} & 24 & Combinación de tipo de horario y categoría \\
\texttt{dim\_\allowbreak{}estado} & 12 & Estado original, estado analítico (R01) y grupo de estado (R04) \\
\texttt{dim\_\allowbreak{}programa} & 175 & Programa normalizado, con el miembro «No informado» \\
\texttt{dim\_\allowbreak{}tipo\_\allowbreak{}encuesta} & 4 & Cuatro tipos: grupo de fuente y etapa (previa o posterior) \\
\texttt{dim\_\allowbreak{}pregunta} & 22 & Texto completo de cada pregunta, con marcas de texto libre y de calificación del tutor \\
\texttt{hecho\_\allowbreak{}reserva} & 86.562 & Un evento de reserva \\
\texttt{hecho\_\allowbreak{}encuesta} & 60.224 & Una encuesta incluida (R02) \\
\texttt{hecho\_\allowbreak{}respuesta} & 384.792 & Una respuesta no vacía a una pregunta \\
\bottomrule
\end{longtable}
\endgroup


\textbf{Conciliación entre Silver y Gold}

\begingroup\footnotesize\setlength{\tabcolsep}{4pt}\renewcommand{\arraystretch}{1.15}
\begin{longtable}{>{\raggedright\arraybackslash}p{\dimexpr 0.262\textwidth-2\tabcolsep\relax}>{\raggedleft\arraybackslash}p{\dimexpr 0.086\textwidth-2\tabcolsep\relax}>{\raggedleft\arraybackslash}p{\dimexpr 0.112\textwidth-2\tabcolsep\relax}>{\raggedleft\arraybackslash}p{\dimexpr 0.112\textwidth-2\tabcolsep\relax}>{\raggedleft\arraybackslash}p{\dimexpr 0.127\textwidth-2\tabcolsep\relax}>{\raggedleft\arraybackslash}p{\dimexpr 0.120\textwidth-2\tabcolsep\relax}>{\raggedright\arraybackslash}p{\dimexpr 0.112\textwidth-2\tabcolsep\relax}}
\toprule
\textbf{Fuente} & \textbf{Filas Silver} & \textbf{Encuestas Gold} & \textbf{Excluidas R02} & \textbf{Estado desconocido pendiente} & \textbf{Respuestas no vacías} & \textbf{Resultado} \\
\midrule
\endhead
Encuesta de satisfacción Express & 6.759 & 6.706 & 53 & 0 & 55.720 & OK \\
Encuesta de satisfacción Normal & 16.379 & 16.268 & 111 & 0 & 180.072 & OK \\
Encuesta previa Express & 13.815 & 13.815 & 0 & 0 & 55.260 & OK \\
Encuesta previa Normal IP & 23.437 & 23.435 & 2 & 0 & 93.740 & OK \\
\bottomrule
\end{longtable}
\endgroup


Esta tabla y los controles siguientes corresponden a la carga anterior a la regla R03. Con R03, el programa de carga también excluye las reservas que no son de tutorías y las encuestas de esas reservas, agrega la columna \texttt{excluidas\_\allowbreak{}R03} a la conciliación y los controles \texttt{reservas\_\allowbreak{}excluidas\_\allowbreak{}R03} y \texttt{encuestas\_\allowbreak{}excluidas\_\allowbreak{}R03}, de modo que las filas de Silver igualan las de Gold más las exclusiones R02 y R03. Controles de la carga anterior: reservas Silver 86.562 y Gold 86.562; atenciones R01 42.646; encuestas excluidas por R02 166; claves foráneas OK; unicidad de granularidades OK; conciliación OK. No se exige igualdad entre el número de respuestas y el de encuestas porque sus granularidades son distintas. Las encuestas con campos vacíos se conservan a nivel de encuesta pero no generan respuestas inventadas.

\textbf{Atributos de las tablas de hechos}

\begingroup\small\setlength{\tabcolsep}{4pt}\renewcommand{\arraystretch}{1.15}
\begin{longtable}{>{\raggedright\arraybackslash}p{\dimexpr 0.105\textwidth-2\tabcolsep\relax}>{\raggedright\arraybackslash}p{\dimexpr 0.156\textwidth-2\tabcolsep\relax}>{\raggedright\arraybackslash}p{\dimexpr 0.335\textwidth-2\tabcolsep\relax}>{\raggedright\arraybackslash}p{\dimexpr 0.335\textwidth-2\tabcolsep\relax}}
\toprule
\textbf{Tabla} & \textbf{Clave primaria} & \textbf{Claves foráneas} & \textbf{Medidas y atributos} \\
\midrule
\endhead
\texttt{hecho\_\allowbreak{}reserva} & \texttt{id\_\allowbreak{}reserva} & fecha de inicio, fin y llegada; hora de inicio; servicio; modalidad; estado; programa & \texttt{evento\_\allowbreak{}reserva} (siempre 1), \texttt{duracion\_\allowbreak{}programada\_\allowbreak{}minutos} ($\geq$ 0), \texttt{es\_\allowbreak{}prioritaria}, fechas originales, archivo, hoja y fila de origen, banderas de calidad \\
\texttt{hecho\_\allowbreak{}encuesta} & \texttt{id\_\allowbreak{}encuesta} & \texttt{id\_\allowbreak{}reserva}, tipo de encuesta, fecha y hora de inicio de la reserva, servicio, modalidad, estado, programa & \texttt{respuesta\_\allowbreak{}encuesta} (siempre 1), \texttt{calificacion\_\allowbreak{}ayuda\_\allowbreak{}tutor} (1 a 5, puede estar vacía), origen, banderas, diferencias con la reserva; restricción única (\texttt{id\_\allowbreak{}reserva}, tipo de encuesta) \\
\texttt{hecho\_\allowbreak{}respuesta} & \texttt{id\_\allowbreak{}encuesta} y \texttt{sk\_\allowbreak{}pregunta} & encuesta, pregunta & \texttt{valor\_\allowbreak{}original} (no vacío), \texttt{valor\_\allowbreak{}numerico} (solo la calificación 1 a 5), \texttt{respuesta\_\allowbreak{}item} (siempre 1) \\
\bottomrule
\end{longtable}
\endgroup


Las referencias de fecha en las encuestas describen la reserva asociada, no el momento en que se contestó. Los conteos son aditivos dentro de su granularidad. Los minutos programados son aditivos como tiempo reservado, pero pueden superponerse entre recursos o usuarios y no equivalen a horas trabajadas ni a ocupación. La calificación no se suma: se calcula su media y mediana con las respuestas no vacías mostrando el tamaño del grupo. Los porcentajes se recalculan desde numerador y denominador. Las claves sustitutas son secuenciales y pueden cambiar si cambia el conjunto de miembros; no deben usarse como identidad estable fuera de una reconstrucción completa.

\subsection*{C.13 Consultas SQL de los análisis}

Se ejecutan sobre la base SQLite de la capa Gold. Gold solo contiene reservas de tutorías (R03) y encuestas válidas (R02). El período académico es un atributo de \texttt{dim\_\allowbreak{}fecha}, y se une por la fecha de inicio de la reserva.

\begingroup\footnotesize\begin{verbatim}
-- Ejecutar sobre data/oro/bookeau.sqlite3. Gold solo tiene reservas de tutorías (R03) y encuestas 
    válidas (R02).
-- El período académico es un atributo de dim_fecha (se une por la fecha de inicio de la reserva).
-- Análisis 1: eventos registrados y atención R01, sin asumir capacidad disponible.
SELECT p.periodo_original, s.servicio_original, m.tipo_horario_original,
       COUNT(*) AS eventos_registrados,
       SUM(CASE WHEN e.estado_analitico='Atendida' THEN 1 ELSE 0 END) AS atenciones_registradas,
       SUM(CASE WHEN e.estado_original='No asistió' THEN 1 ELSE 0 END) AS inasistencias_registradas
FROM hecho_reserva h
JOIN dim_fecha p ON p.sk_fecha=h.sk_fecha_inicio
JOIN dim_servicio s USING(sk_servicio)
JOIN dim_modalidad m USING(sk_modalidad)
JOIN dim_estado e USING(sk_estado)
GROUP BY p.periodo_original,s.servicio_original,m.tipo_horario_original;

-- Análisis 2: calificación de encuestas posteriores incluidas, con denominador del ítem.
SELECT p.periodo_original, s.servicio_original, m.tipo_horario_original,
       COUNT(*) AS encuestas_posteriores_incluidas,
       COUNT(h.calificacion_ayuda_tutor) AS n_calificaciones,
       AVG(h.calificacion_ayuda_tutor) AS media_descriptiva,
       100.0*COUNT(h.calificacion_ayuda_tutor)/COUNT(*) AS completitud_item_pct
FROM hecho_encuesta h
JOIN dim_tipo_encuesta t USING(sk_tipo_encuesta)
JOIN dim_fecha p ON p.sk_fecha=h.sk_fecha_inicio
JOIN dim_servicio s USING(sk_servicio)
JOIN dim_modalidad m USING(sk_modalidad)
WHERE t.etapa='posterior'
GROUP BY p.periodo_original,s.servicio_original,m.tipo_horario_original;

-- Texto válido disponible para etapas futuras; no mezclar su grano con reservas.
SELECT q.texto_pregunta,t.grupo_fuente,COUNT(*) AS respuestas_textuales
FROM hecho_respuesta r
JOIN dim_pregunta q USING(sk_pregunta)
JOIN hecho_encuesta e USING(id_encuesta)
JOIN dim_tipo_encuesta t USING(sk_tipo_encuesta)
WHERE q.es_texto_libre=1
GROUP BY q.texto_pregunta,t.grupo_fuente;

-- Análisis 1 (R04): tasa de inasistencia y proporción en cola por modalidad.
-- dim_estado.grupo_estado materializa R04: inasistencia = No asistió / (Atendida + No asistió).
SELECT m.tipo_horario_original AS modalidad,
       COUNT(*) AS eventos,
       ROUND(100.0*SUM(e.grupo_estado='No asistió')
             /SUM(e.grupo_estado IN ('Atendida','No asistió')),1) AS inasistencia_pct,
       ROUND(100.0*SUM(e.grupo_estado='Cola')/COUNT(*),1) AS cola_pct
FROM hecho_reserva h
JOIN dim_modalidad m USING(sk_modalidad)
JOIN dim_estado e USING(sk_estado)
GROUP BY m.tipo_horario_original
ORDER BY eventos DESC;

-- Análisis 1 (R04): inasistencia por día de la semana y franja horaria (jerarquías de Fecha y 
    Hora).
SELECT d.dia_semana_iso, d.nombre_dia, t.franja,
       SUM(e.grupo_estado IN ('Atendida','No asistió')) AS citas_programadas,
       ROUND(100.0*SUM(e.grupo_estado='No asistió')
             /SUM(e.grupo_estado IN ('Atendida','No asistió')),1) AS inasistencia_pct
FROM hecho_reserva h
JOIN dim_fecha d ON d.sk_fecha=h.sk_fecha_inicio
JOIN dim_hora t ON t.sk_hora=h.sk_hora_inicio
JOIN dim_estado e USING(sk_estado)
GROUP BY d.dia_semana_iso, d.nombre_dia, t.franja
ORDER BY d.dia_semana_iso, t.franja;

-- Análisis 1: eventos y cola por tipo de período (atributos de período en dim_fecha).
SELECT p.tipo_periodo, COUNT(*) AS eventos,
       ROUND(100.0*SUM(e.grupo_estado='Cola')/COUNT(*),1) AS cola_pct
FROM hecho_reserva h
JOIN dim_fecha p ON p.sk_fecha=h.sk_fecha_inicio
JOIN dim_estado e USING(sk_estado)
GROUP BY p.tipo_periodo;

-- Análisis 2: proporción de calificaciones 1-3 por período (encuestas posteriores válidas).
SELECT p.periodo_original,
       COUNT(h.calificacion_ayuda_tutor) AS n_calificaciones,
       ROUND(100.0*SUM(h.calificacion_ayuda_tutor<=3)/COUNT(h.calificacion_ayuda_tutor),1) AS 
    bajas_pct
FROM hecho_encuesta h
JOIN dim_tipo_encuesta t USING(sk_tipo_encuesta)
JOIN dim_fecha p ON p.sk_fecha=h.sk_fecha_inicio
WHERE t.etapa='posterior'
GROUP BY p.periodo_original
HAVING COUNT(h.calificacion_ayuda_tutor)>=30;
\end{verbatim}\endgroup

\textbf{Definición de las tablas (DDL).}

\begingroup\footnotesize\begin{verbatim}
PRAGMA foreign_keys = ON;
CREATE TABLE dim_fecha(sk_fecha INTEGER PRIMARY KEY, fecha TEXT UNIQUE NOT NULL, anio INTEGER, mes 
    INTEGER, dia INTEGER, dia_semana_iso INTEGER CHECK(dia_semana_iso BETWEEN 1 AND 7), nombre_dia 
    TEXT NOT NULL, nombre_mes TEXT NOT NULL, es_fin_de_semana INTEGER CHECK(es_fin_de_semana IN 
    (0,1)), periodo_original TEXT NOT NULL, anio_periodo INTEGER, sufijo_periodo TEXT, tipo_periodo 
    TEXT NOT NULL);
CREATE TABLE dim_hora(sk_hora INTEGER PRIMARY KEY, hora INTEGER, minuto INTEGER, etiqueta TEXT, 
    franja TEXT NOT NULL CHECK(franja IN ('Madrugada','Mañana','Tarde','Noche')));
CREATE TABLE dim_servicio(sk_servicio INTEGER PRIMARY KEY, codigo_servicio TEXT, servicio_original 
    TEXT, UNIQUE(codigo_servicio, servicio_original));
CREATE TABLE dim_modalidad(sk_modalidad INTEGER PRIMARY KEY, tipo_horario_original TEXT, 
    categoria_original TEXT, UNIQUE(tipo_horario_original,categoria_original));
CREATE TABLE dim_estado(sk_estado INTEGER PRIMARY KEY, estado_original TEXT UNIQUE NOT NULL, 
    estado_analitico TEXT NOT NULL, grupo_estado TEXT NOT NULL CHECK(grupo_estado IN ('Atendida','No
     asistió','Cola','Cancelada','Abierta')));
CREATE TABLE dim_programa(sk_programa INTEGER PRIMARY KEY, programa_normalizado TEXT UNIQUE NOT 
    NULL);
CREATE TABLE dim_tipo_encuesta(sk_tipo_encuesta INTEGER PRIMARY KEY, grupo_fuente TEXT UNIQUE NOT 
    NULL, etapa TEXT NOT NULL);
CREATE TABLE dim_pregunta(sk_pregunta INTEGER PRIMARY KEY, columna_plata TEXT UNIQUE NOT NULL, 
    texto_pregunta TEXT NOT NULL, es_texto_libre INTEGER CHECK(es_texto_libre IN (0,1)), 
    es_calificacion_tutor INTEGER CHECK(es_calificacion_tutor IN (0,1)));
CREATE TABLE hecho_reserva(
 id_reserva TEXT PRIMARY KEY, sk_fecha_inicio INTEGER REFERENCES dim_fecha, sk_fecha_fin INTEGER 
    REFERENCES dim_fecha,
 sk_fecha_llegada INTEGER REFERENCES dim_fecha, sk_hora_inicio INTEGER NOT NULL REFERENCES dim_hora,
 sk_servicio INTEGER NOT NULL REFERENCES dim_servicio,
 sk_modalidad INTEGER NOT NULL REFERENCES dim_modalidad, sk_estado INTEGER NOT NULL REFERENCES 
    dim_estado,
 sk_programa INTEGER NOT NULL REFERENCES dim_programa,
 fecha_inicio TEXT NOT NULL, fecha_fin TEXT NOT NULL, fecha_llegada TEXT,
 duracion_programada_minutos REAL CHECK(duracion_programada_minutos>=0), es_prioritaria INTEGER 
    CHECK(es_prioritaria IN (0,1)),
 evento_reserva INTEGER NOT NULL DEFAULT 1 CHECK(evento_reserva=1),
 archivo_origen TEXT NOT NULL, hoja_origen TEXT NOT NULL, fila_excel INTEGER NOT NULL, 
    banderas_calidad_json TEXT NOT NULL);
CREATE TABLE hecho_encuesta(
 id_encuesta TEXT PRIMARY KEY, id_reserva TEXT NOT NULL REFERENCES hecho_reserva,
 sk_tipo_encuesta INTEGER NOT NULL REFERENCES dim_tipo_encuesta,
 sk_fecha_inicio INTEGER REFERENCES dim_fecha, sk_hora_inicio INTEGER NOT NULL REFERENCES dim_hora,
 sk_servicio INTEGER NOT NULL REFERENCES dim_servicio,
 sk_modalidad INTEGER NOT NULL REFERENCES dim_modalidad, sk_estado INTEGER NOT NULL REFERENCES 
    dim_estado,
 sk_programa INTEGER NOT NULL REFERENCES dim_programa,
 calificacion_ayuda_tutor INTEGER CHECK(calificacion_ayuda_tutor BETWEEN 1 AND 5),
 respuesta_encuesta INTEGER NOT NULL DEFAULT 1 CHECK(respuesta_encuesta=1),
 archivo_origen TEXT NOT NULL, hoja_origen TEXT NOT NULL, fila_excel INTEGER NOT NULL,
 banderas_calidad_json TEXT NOT NULL, diferencias_con_reserva_json TEXT NOT NULL,
 UNIQUE(id_reserva,sk_tipo_encuesta));
CREATE TABLE hecho_respuesta(
 id_encuesta TEXT NOT NULL REFERENCES hecho_encuesta,
 sk_pregunta INTEGER NOT NULL REFERENCES dim_pregunta,
 valor_original TEXT NOT NULL CHECK(length(trim(valor_original))>0),
 valor_numerico REAL, respuesta_item INTEGER NOT NULL DEFAULT 1 CHECK(respuesta_item=1),
 PRIMARY KEY(id_encuesta,sk_pregunta));
\end{verbatim}\endgroup

\subsection*{C.14 Resultados adicionales de los análisis}

\textbf{Indicadores globales}

\begingroup\small\setlength{\tabcolsep}{4pt}\renewcommand{\arraystretch}{1.15}
\begin{longtable}{>{\raggedright\arraybackslash}p{\dimexpr 0.736\textwidth-2\tabcolsep\relax}>{\raggedleft\arraybackslash}p{\dimexpr 0.194\textwidth-2\tabcolsep\relax}}
\toprule
\textbf{Indicador} & \textbf{Valor} \\
\midrule
\endhead
Eventos de reserva & 86.562 \\
Atenciones (R01) & 42.646 \\
Eventos «No asistió» & 5.086 \\
Encuestas posteriores incluidas (R02) & 22.974 \\
Calificaciones utilizables & 16.171 \\
Media descriptiva de la calificación & 4,74 \\
Tasa de inasistencia (R04) & 10,7\% \\
Proporción en cola (R04) & 17,1\% \\
Proporción de calificaciones 1 a 3 & 5,7\% \\
\bottomrule
\end{longtable}
\endgroup


\textbf{Satisfacción por modalidad original} (todas las modalidades presentes, incluidas las obsoletas)

\begingroup\small\setlength{\tabcolsep}{4pt}\renewcommand{\arraystretch}{1.15}
\begin{longtable}{>{\raggedright\arraybackslash}p{\dimexpr 0.332\textwidth-2\tabcolsep\relax}>{\raggedleft\arraybackslash}p{\dimexpr 0.178\textwidth-2\tabcolsep\relax}>{\raggedleft\arraybackslash}p{\dimexpr 0.216\textwidth-2\tabcolsep\relax}>{\raggedleft\arraybackslash}p{\dimexpr 0.204\textwidth-2\tabcolsep\relax}}
\toprule
\textbf{Modalidad original} & \textbf{Encuestas incluidas} & \textbf{Con calificación} & \textbf{Media descriptiva} \\
\midrule
\endhead
Express & 6.646 & 4.878 & 4,6982 \\
Normal & 14.731 & 9.722 & 4,7725 \\
Normal pico & 1.450 & 1.450 & 4,6241 \\
Deprecated, Grupal & 60 & 60 & 4,6000 \\
Deprecated, Torre Séneca & 87 & 61 & 4,7049 \\
\bottomrule
\end{longtable}
\endgroup


Faltan 6.803 calificaciones y no se les asigna cero; la completitud del ítem entre las encuestas incluidas es de 70,39\%. La mediana global es 5. Las medias agregadas no demuestran superioridad de una modalidad: hay que comparar por período, servicio, muestra y formulario, y la media supone distancias iguales entre los puntos de la escala. La ausencia de respuestas y los cambios históricos pueden sesgar la comparación, y no se infiere una mejora causal del aprendizaje.

\textbf{Servicios y modalidades con más eventos (etiquetas originales).} Servicios: IP 53.272, Deprecated APO 1 21.136, EDA 6.127. Modalidades: Normal 49.472, Express 25.754, Normal pico 8.668. Hora de inicio programado con más eventos: 11:00 a 11:59 (13.699), 14:00 a 14:59 (11.253) y 09:00 a 09:59 (10.217). Estos conteos reflejan agenda reservada y no demuestran insuficiencia de oferta. Los totales incluyen modalidades históricas, salas y entrevistas, que el tablero permite separar.

\textbf{Texto disponible en Gold} (respuestas textuales de encuestas incluidas)

\begingroup\small\setlength{\tabcolsep}{4pt}\renewcommand{\arraystretch}{1.15}
\begin{longtable}{>{\raggedright\arraybackslash}p{\dimexpr 0.494\textwidth-2\tabcolsep\relax}>{\raggedright\arraybackslash}p{\dimexpr 0.299\textwidth-2\tabcolsep\relax}>{\raggedleft\arraybackslash}p{\dimexpr 0.136\textwidth-2\tabcolsep\relax}}
\toprule
\textbf{Pregunta} & \textbf{Encuesta} & \textbf{Respuestas} \\
\midrule
\endhead
Comente aspectos positivos y negativos del tutor y la tutoría. Recuerde que las respuestas van a ser anónimas. & Encuesta de satisfacción Express & 5.403 \\
Comente aspectos positivos y negativos del tutor y la tutoría. Recuerde que las respuestas van a ser anónimas. & Encuesta de satisfacción Normal & 12.426 \\
Escriba en este recuadro cualquier observación o comentario relacionado al curso y sus contenidos. & Encuesta de satisfacción Express & 561 \\
Escriba en este recuadro cualquier observación o comentario relacionado al curso y sus contenidos. & Encuesta de satisfacción Normal & 1.317 \\
Indique cuál fue la herramienta usada durante la tutoría (si sabe cuál es) y el uso que tuvo la misma en la tutoría. & Encuesta de satisfacción Express & 4 \\
Indique cuál fue la herramienta usada durante la tutoría (si sabe cuál es) y el uso que tuvo la misma en la tutoría. & Encuesta de satisfacción Normal & 22 \\
Sugerencias, reclamos y observaciones & Encuesta de satisfacción Express & 1.703 \\
Sugerencias, reclamos y observaciones & Encuesta de satisfacción Normal & 4.005 \\
\bottomrule
\end{longtable}
\endgroup


\textbf{Archivos que consume el tablero.} Agregados por período, servicio, modalidad, estado, día y hora, y distribución de calificaciones. Contienen solo información de grupos, sin nombres ni correos.

\section*{Anexo D · Extensión exploratoria: comentarios por monitor}

Esta extensión no es un requisito de la primera entrega: la guía reserva el uso del texto para las siguientes. Se incluye como avance. Parte de \texttt{hecho\_\allowbreak{}respuesta} y de la encuesta de Gold, y toma el prestador de Silver.

\textbf{Objetivo.} Explorar qué aspectos se valoran o cuestionan en los comentarios sobre tutor y tutoría, y cómo se distribuyen por monitor, período, servicio y modalidad. Se usa solo la pregunta que pide aspectos positivos y negativos del tutor y de la tutoría; los comentarios sobre el curso no se atribuyen al monitor.

\textbf{Corpus e identificación.} Hay 17.829 comentarios con contenido de encuestas posteriores incluidas por R02. En 17.534 coincide el nombre del prestador entre la encuesta y la reserva, después de normalizar mayúsculas, Unicode y espacios. Resultan 424 etiquetas de monitor con comentarios atribuidos; son etiquetas normalizadas, no 424 identidades certificadas. Los 295 comentarios restantes quedan sin monitor confirmado: 291 con nombres distintos entre las fuentes y 4 sin monitor en ninguna. La pregunta habla de anonimato, pero el paquete permite vincular ID y prestador; este análisis usa esa vinculación solo en el entorno local del equipo.

\begingroup\small\setlength{\tabcolsep}{4pt}\renewcommand{\arraystretch}{1.15}
\begin{longtable}{>{\raggedright\arraybackslash}p{\dimexpr 0.315\textwidth-2\tabcolsep\relax}>{\raggedright\arraybackslash}p{\dimexpr 0.615\textwidth-2\tabcolsep\relax}}
\toprule
\textbf{Categoría} & \textbf{Interpretación de la regla} \\
\midrule
\endhead
Positivo & Hay señales favorables reconocidas y ninguna desfavorable \\
Negativo & Hay señales desfavorables reconocidas y ninguna favorable \\
Mixto & Se reconocen señales de ambos tipos \\
Sin información suficiente & Solo símbolos o respuestas genéricas (nada, ninguno, N/A) \\
Por revisar & Texto con contenido pero sin señal reconocida por las reglas \\
\bottomrule
\end{longtable}
\endgroup


Cada registro conserva el texto original, las reglas activadas, la etiqueta automática, la etiqueta final, el origen de la etiqueta y los temas. Las negaciones se revisan en una ventana corta y el contexto se limita por signos de puntuación; el método no resuelve ironía, errores de escritura ni negaciones complejas. Temas detectados por menciones léxicas, que pueden solaparse: claridad, resolución de dudas, trato, tiempo, conocimiento, puntualidad, aprendizaje y conectividad o modalidad.

\textbf{Distribución automática actual (versión 2)}

\begingroup\small\setlength{\tabcolsep}{4pt}\renewcommand{\arraystretch}{1.15}
\begin{longtable}{>{\raggedright\arraybackslash}p{\dimexpr 0.591\textwidth-2\tabcolsep\relax}>{\raggedleft\arraybackslash}p{\dimexpr 0.339\textwidth-2\tabcolsep\relax}}
\toprule
\textbf{Etiqueta} & \textbf{Comentarios} \\
\midrule
\endhead
Positivo & 12.299 \\
Sin información suficiente & 2.569 \\
Por revisar & 2.273 \\
Mixto & 361 \\
Negativo & 327 \\
\bottomrule
\end{longtable}
\endgroup


Son predicciones provisionales. Las 17.579 etiquetas sin revisión individual no justifican un ranking ni tasas de desempeño por monitor.

\textbf{Evaluación del clasificador.} La revisión de referencia la hizo una IA y no personas; el mismo asistente desarrolló las reglas y produjo la referencia, de modo que no hay independencia entre desarrollador y evaluador. Procedimiento: se conservaron las predicciones originales; se revisaron 150 comentarios de desarrollo (30 por categoría automática) leídos sin predicciones ni nombres de monitor; se ajustaron las reglas; se congelaron las predicciones v2; se seleccionaron otros 100 comentarios (20 por categoría) sin textos vistos en desarrollo; y no se ajustó nada después de leer esa evaluación reservada.

\begingroup\small\setlength{\tabcolsep}{4pt}\renewcommand{\arraystretch}{1.15}
\begin{longtable}{>{\raggedright\arraybackslash}p{\dimexpr 0.297\textwidth-2\tabcolsep\relax}>{\raggedright\arraybackslash}p{\dimexpr 0.336\textwidth-2\tabcolsep\relax}>{\raggedright\arraybackslash}p{\dimexpr 0.297\textwidth-2\tabcolsep\relax}}
\toprule
\textbf{Versión} & \textbf{Desarrollo (150)} & \textbf{Evaluación reservada (100)} \\
\midrule
\endhead
v1 original & 88/150 (58,7\%) & 57/100 (57\%) \\
v2 corregida & 150/150 (100\%) & 62/100 (62\%) \\
\bottomrule
\end{longtable}
\endgroup


El 100\% de desarrollo no demuestra generalización. Resultados de v2 por categoría en la evaluación reservada:

\begingroup\small\setlength{\tabcolsep}{4pt}\renewcommand{\arraystretch}{1.15}
\begin{longtable}{>{\raggedright\arraybackslash}p{\dimexpr 0.266\textwidth-2\tabcolsep\relax}>{\raggedleft\arraybackslash}p{\dimexpr 0.153\textwidth-2\tabcolsep\relax}>{\raggedleft\arraybackslash}p{\dimexpr 0.162\textwidth-2\tabcolsep\relax}>{\raggedleft\arraybackslash}p{\dimexpr 0.134\textwidth-2\tabcolsep\relax}>{\raggedleft\arraybackslash}p{\dimexpr 0.134\textwidth-2\tabcolsep\relax}>{\raggedleft\arraybackslash}p{\dimexpr 0.082\textwidth-2\tabcolsep\relax}}
\toprule
\textbf{Categoría} & \textbf{Referencias IA} & \textbf{Predicciones} & \textbf{Precisión} & \textbf{Recobrado} & \textbf{F1} \\
\midrule
\endhead
Positivo & 46 & 31 & 80,6\% & 54,3\% & 0,649 \\
Negativo & 21 & 15 & 73,3\% & 52,4\% & 0,611 \\
Mixto & 13 & 16 & 43,8\% & 53,8\% & 0,483 \\
Sin información suficiente & 18 & 20 & 90,0\% & 100\% & 0,947 \\
Por revisar & 2 & 18 & 5,6\% & 50,0\% & 0,100 \\
\bottomrule
\end{longtable}
\endgroup


El F1 macro de las cinco categorías es 0,558. La precisión es el número de coincidencias de una categoría dividido por sus predicciones; el recobrado es ese número dividido por las referencias IA de la categoría. «Por revisar» es una abstención operativa y no una polaridad. El 62\% no estima la exactitud sobre los 17.829 comentarios, ni los porcentajes de la muestra estiman la prevalencia por monitor. Conclusión: \textbf{el clasificador de reglas todavía no es suficientemente fiable para evaluar o comparar monitores.}

\textbf{Errores observados}

\begingroup\small\setlength{\tabcolsep}{4pt}\renewcommand{\arraystretch}{1.15}
\begin{longtable}{>{\raggedright\arraybackslash}p{\dimexpr 0.373\textwidth-2\tabcolsep\relax}>{\raggedright\arraybackslash}p{\dimexpr 0.107\textwidth-2\tabcolsep\relax}>{\raggedright\arraybackslash}p{\dimexpr 0.093\textwidth-2\tabcolsep\relax}>{\raggedright\arraybackslash}p{\dimexpr 0.357\textwidth-2\tabcolsep\relax}}
\toprule
\textbf{Comentario de evaluación} & \textbf{Predicción v2} & \textbf{Revisión IA} & \textbf{Problema} \\
\midrule
\endhead
Todo perfecto & Por revisar & Positivo & Falta vocabulario favorable \\
El tutor fue un poco grosero y agresivo hacia mi persona & Por revisar & Negativo & «poco» suprime indebidamente una crítica \\
A pesar del poco tiempo fue capaz de aclararme dudas & Negativo & Positivo & La restricción de tiempo domina el beneficio expresado \\
Me permitió ver los errores que tenía y como corregirlos. & Por revisar & Positivo & No reconoce una valoración favorable implícita \\
\bottomrule
\end{longtable}
\endgroup


Otros desacuerdos mezclan elogios con problemas de conexión, sugerencias de mejora o críticas al trato, y requieren decidir si el comentario critica al tutor, a la sesión o a una restricción externa. Si se usa esta evaluación para desarrollar otra versión, deja de ser reservada y hará falta otro conjunto de evaluación.

\textbf{Tablero por monitor y revisión humana.} El explorador local permite filtrar por monitor, período, servicio, modalidad y clasificación, buscar dentro del comentario, consultar temas y corregir etiquetas durante la sesión. Muestra por defecto los comentarios revisados y expone conteos, no porcentajes, porque la muestra revisada es estratificada. Se preparó una muestra de 150 comentarios (30 por categoría automática, con semilla fija) para etiquetado humano; sus etiquetas siguen vacías. La coincidencia de esa muestra no equivaldría a precisión poblacional.

\section*{Anexo E · Reglas pendientes de confirmar y estado del proyecto}

\textbf{Reglas que requieren confirmación de la fuente}

\begin{enumerate}[leftmargin=1.5em,itemsep=2pt]
\item La diferencia técnica entre «Finalizada» y «Realizada» y la clasificación de los demás estados. Hoy solo se adopta R01 como equivalencia de trabajo aportada por la coordinación.
\item La población elegible para las tasas de asistencia y cancelación, el tratamiento de la cola y el momento en que se registra la llegada.
\item El criterio de encuesta «Inválida», la población invitada y si la encuesta es obligatoria.
\item Los catálogos históricos de servicios y modalidades, el significado de los períodos y la zona horaria de las marcas de tiempo.
\item Los cuestionarios y escalas por período, y si la exportación usa etiquetas actuales para respuestas históricas.
\item La fuente de referencia y el momento de actualización de los atributos de reserva repetidos en las encuestas.
\end{enumerate}

Mientras se confirman, pueden presentarse el volumen de eventos, la agenda programada, las distribuciones de etiquetas y las calificaciones del subconjunto documentado. Los roles destinatarios de los análisis son propuestos: no se ha confirmado el organigrama ni se ha entrevistado al responsable de la fuente.

\textbf{Limitaciones de los datos.} No hay oferta ni capacidad general de franjas, fin efectivo de la sesión, respuestas del tutor como fuente separada, ni versiones certificadas del cuestionario. No se publican ocupación, cupos libres ni demanda insatisfecha. Tampoco hay notas académicas ni medidas objetivas de aprendizaje, costos o ingresos, de modo que los datos no permiten medir impacto causal en el aprendizaje, rentabilidad ni uso de capacidad. Para incorporar la oferta harían falta registros de franjas ofrecidas, capacidad, recursos y cambios de disponibilidad, con fechas y claves compatibles; su granularidad sería distinta de la de reserva y el modelo no fabrica esos eventos.

\textbf{Estado de las tareas de la entrega}

\begingroup\small\setlength{\tabcolsep}{4pt}\renewcommand{\arraystretch}{1.15}
\begin{longtable}{>{\raggedright\arraybackslash}p{\dimexpr 0.763\textwidth-2\tabcolsep\relax}>{\raggedright\arraybackslash}p{\dimexpr 0.167\textwidth-2\tabcolsep\relax}}
\toprule
\textbf{Tarea} & \textbf{Estado} \\
\midrule
\endhead
Espacio de trabajo, PDF del enunciado y ZIP original de Bookeau & Hecho \\
Extracción de las fuentes e inventario por archivo, hoja, período y encabezados & Hecho \\
Caracterización completa, perfilamiento y verificación de ID y vínculos & Hecho \\
Definición de los dos análisis con objetivos, medidas, preguntas y denominadores & Hecho \\
Evidencia de calidad y agregados descriptivos & Hecho \\
Limpieza Bronze a Silver con trazabilidad, banderas y conciliación & Hecho \\
Ecosistema propuesto y justificado & Hecho \\
Modelo: granularidad, hechos, dimensiones y criterios de calidad & Hecho \\
Transformación Silver a Gold con conciliación & Hecho \\
Tablero local revisado visualmente y respuestas documentadas & Hecho \\
Cuadernos ejecutados e informe con tablas, diagramas y capturas & Hecho \\
Confirmar integrantes, contexto organizacional y roles destinatarios & Pendiente \\
Revisar los casos marcados y completar reglas operativas de limpieza & Pendiente \\
Confirmar con la fuente reglas operativas, catálogos y formularios históricos & Pendiente \\
Validar con la coordinación la interpretación de los resultados & Pendiente \\
Completar integrantes, contribuciones y distribución de los 100 puntos, y revisión final & Pendiente \\
\bottomrule
\end{longtable}
\endgroup


\end{document}
